% Môn: Vật lý
% Lớp: 10
% Chương: Chương III. Động lực học
% Bài: L10C3 Bài 15. Định luật 2 Newton
% Số câu: 162
% App đọc file này trực tiếp — sửa rồi Commit trên GitHub, bấm Đồng bộ GitHub .tex

% L10C3 Bài 15. Định luật 2 Newton

% ===== Câu 1 =====
% ID: AIQ_L10C3_FULL_0325
% Mức: NB
\dangbt{Nhận biết nội dung và biểu thức định luật II Newton}
\begin{ex}
% ID: AIQ_L10C3_FULL_0325
% Mức: NB
Vật khối lượng $2\,\text{kg}$ chịu hợp lực $2\,\text{N}$. Tính gia tốc (kết quả theo đơn vị m/s^2).
\choice
{\True 1}
{2}
{0}
{2}
\loigiai{
Áp dụng hệ thức phù hợp của bài học, thay số được $1\,\text{m}$ /s^2.
}
\end{ex}

% ===== Câu 2 =====
% ID: AIQ_L10C3_FULL_0326
% Mức: NB
\begin{ex}
% ID: AIQ_L10C3_FULL_0326
% Mức: NB
Vật khối lượng $3\,\text{kg}$ chịu hợp lực $6\,\text{N}$. Tính gia tốc (kết quả theo đơn vị m/s^2).
\choice
{3}
{\True 2}
{1}
{4}
\loigiai{
Áp dụng hệ thức phù hợp của bài học, thay số được $2\,\text{m}$ /s^2.
}
\end{ex}

% ===== Câu 3 =====
% ID: AIQ_L10C3_FULL_0327
% Mức: NB
\begin{ex}
% ID: AIQ_L10C3_FULL_0327
% Mức: NB
Vật khối lượng $4\,\text{kg}$ chịu hợp lực $12\,\text{N}$. Tính gia tốc (kết quả theo đơn vị m/s^2).
\choice
{2}
{4}
{\True 3}
{6}
\loigiai{
Áp dụng hệ thức phù hợp của bài học, thay số được $3\,\text{m}$ /s^2.
}
\end{ex}

% ===== Câu 4 =====
% ID: AIQ_L10C3_FULL_0328
% Mức: NB
\begin{ex}
% ID: AIQ_L10C3_FULL_0328
% Mức: NB
Vật khối lượng $5\,\text{kg}$ chịu hợp lực $20\,\text{N}$. Tính gia tốc (kết quả theo đơn vị m/s^2).
\choice
{8}
{5}
{3}
{\True 4}
\loigiai{
Áp dụng hệ thức phù hợp của bài học, thay số được $4\,\text{m}$ /s^2.
}
\end{ex}

% ===== Câu 5 =====
% ID: AIQ_L10C3_FULL_0329
% Mức: TH
\begin{ex}
% ID: AIQ_L10C3_FULL_0329
% Mức: TH
Vật khối lượng $6\,\text{kg}$ chịu hợp lực $6\,\text{N}$. Tính gia tốc (kết quả theo đơn vị m/s^2).
\choice
{\True 1}
{2}
{0}
{2}
\loigiai{
Áp dụng hệ thức phù hợp của bài học, thay số được $1\,\text{m}$ /s^2.
}
\end{ex}

% ===== Câu 6 =====
% ID: AIQ_L10C3_FULL_0330
% Mức: TH
\begin{ex}
% ID: AIQ_L10C3_FULL_0330
% Mức: TH
Vật khối lượng $2\,\text{kg}$ chịu hợp lực $4\,\text{N}$. Tính gia tốc (kết quả theo đơn vị m/s^2).
\choice
{3}
{\True 2}
{1}
{4}
\loigiai{
Áp dụng hệ thức phù hợp của bài học, thay số được $2\,\text{m}$ /s^2.
}
\end{ex}

% ===== Câu 7 =====
% ID: AIQ_L10C3_FULL_0331
% Mức: TH
\begin{ex}
% ID: AIQ_L10C3_FULL_0331
% Mức: TH
Vật khối lượng $3\,\text{kg}$ chịu hợp lực $9\,\text{N}$. Tính gia tốc (kết quả theo đơn vị m/s^2).
\choice
{2}
{4}
{\True 3}
{6}
\loigiai{
Áp dụng hệ thức phù hợp của bài học, thay số được $3\,\text{m}$ /s^2.
}
\end{ex}

% ===== Câu 8 =====
% ID: AIQ_L10C3_FULL_0332
% Mức: TH
\begin{ex}
% ID: AIQ_L10C3_FULL_0332
% Mức: TH
Vật khối lượng $4\,\text{kg}$ chịu hợp lực $16\,\text{N}$. Tính gia tốc (kết quả theo đơn vị m/s^2).
\choice
{8}
{5}
{3}
{\True 4}
\loigiai{
Áp dụng hệ thức phù hợp của bài học, thay số được $4\,\text{m}$ /s^2.
}
\end{ex}

% ===== Câu 9 =====
% ID: AIQ_L10C3_FULL_0333
% Mức: VD
\begin{ex}
% ID: AIQ_L10C3_FULL_0333
% Mức: VD
Vật khối lượng $5\,\text{kg}$ chịu hợp lực $5\,\text{N}$. Tính gia tốc (kết quả theo đơn vị m/s^2).
\choice
{\True 1}
{2}
{0}
{2}
\loigiai{
Áp dụng hệ thức phù hợp của bài học, thay số được $1\,\text{m}$ /s^2.
}
\end{ex}

% ===== Câu 10 =====
% ID: AIQ_L10C3_FULL_0334
% Mức: VD
\begin{ex}
% ID: AIQ_L10C3_FULL_0334
% Mức: VD
Vật khối lượng $6\,\text{kg}$ chịu hợp lực $12\,\text{N}$. Tính gia tốc (kết quả theo đơn vị m/s^2).
\choice
{3}
{\True 2}
{1}
{4}
\loigiai{
Áp dụng hệ thức phù hợp của bài học, thay số được $2\,\text{m}$ /s^2.
}
\end{ex}

% ===== Câu 11 =====
% ID: AIQ_L10C3_FULL_0335
% Mức: TH
\begin{ex}
% ID: AIQ_L10C3_FULL_0335
% Mức: TH
Vật khối lượng $2\,\text{kg}$ chịu hợp lực $6\,\text{N}$. Tính gia tốc (kết quả theo đơn vị m/s^2).
\choiceTF
{\True Giá trị cần tìm là $3\,\text{m}$ /s^2.}
{Giá trị cần tìm là $4\,\text{m}$ /s^2.}
{\True Phải xác định đầy đủ các lực hoặc đại lượng liên quan trước khi tính.}
{Có thể bỏ qua điều kiện áp dụng của công thức.}
\loigiai{
\begin{itemchoice}
\item (Đúng) Giá trị cần tìm là $3\,\text{m}$ /s^2. (Vì): kết quả $3\,\text{m}$ /s^2. b) Sai. c) Đúng. d) Sai vì phải kiểm tra điều kiện áp dụng
\item (Sai) Giá trị cần tìm là $4\,\text{m}$ /s^2.
\item (Đúng) Phải xác định đầy đủ các lực hoặc đại lượng liên quan trước khi tính.
\item (Sai) Có thể bỏ qua điều kiện áp dụng của công thức.
\end{itemchoice}
}
\end{ex}

% ===== Câu 12 =====
% ID: AIQ_L10C3_FULL_0336
% Mức: TH
\begin{ex}
% ID: AIQ_L10C3_FULL_0336
% Mức: TH
Vật khối lượng $3\,\text{kg}$ chịu hợp lực $12\,\text{N}$. Tính gia tốc (kết quả theo đơn vị m/s^2).
\choiceTF
{\True Giá trị cần tìm là $4\,\text{m}$ /s^2.}
{Giá trị cần tìm là $5\,\text{m}$ /s^2.}
{\True Phải xác định đầy đủ các lực hoặc đại lượng liên quan trước khi tính.}
{Có thể bỏ qua điều kiện áp dụng của công thức.}
\loigiai{
\begin{itemchoice}
\item (Đúng) Giá trị cần tìm là $4\,\text{m}$ /s^2. (Vì): kết quả $4\,\text{m}$ /s^2. b) Sai. c) Đúng. d) Sai vì phải kiểm tra điều kiện áp dụng
\item (Sai) Giá trị cần tìm là $5\,\text{m}$ /s^2.
\item (Đúng) Phải xác định đầy đủ các lực hoặc đại lượng liên quan trước khi tính.
\item (Sai) Có thể bỏ qua điều kiện áp dụng của công thức.
\end{itemchoice}
}
\end{ex}

% ===== Câu 13 =====
% ID: AIQ_L10C3_FULL_0337
% Mức: VD
\begin{ex}
% ID: AIQ_L10C3_FULL_0337
% Mức: VD
Vật khối lượng $4\,\text{kg}$ chịu hợp lực $4\,\text{N}$. Tính gia tốc (kết quả theo đơn vị m/s^2).
\choiceTF
{\True Giá trị cần tìm là $1\,\text{m}$ /s^2.}
{Giá trị cần tìm là $2\,\text{m}$ /s^2.}
{\True Phải xác định đầy đủ các lực hoặc đại lượng liên quan trước khi tính.}
{Có thể bỏ qua điều kiện áp dụng của công thức.}
\loigiai{
\begin{itemchoice}
\item (Đúng) Giá trị cần tìm là $1\,\text{m}$ /s^2. (Vì): kết quả $1\,\text{m}$ /s^2. b) Sai. c) Đúng. d) Sai vì phải kiểm tra điều kiện áp dụng
\item (Sai) Giá trị cần tìm là $2\,\text{m}$ /s^2.
\item (Đúng) Phải xác định đầy đủ các lực hoặc đại lượng liên quan trước khi tính.
\item (Sai) Có thể bỏ qua điều kiện áp dụng của công thức.
\end{itemchoice}
}
\end{ex}

% ===== Câu 14 =====
% ID: AIQ_L10C3_FULL_0338
% Mức: VD
\begin{ex}
% ID: AIQ_L10C3_FULL_0338
% Mức: VD
Vật khối lượng $5\,\text{kg}$ chịu hợp lực $10\,\text{N}$. Tính gia tốc (kết quả theo đơn vị m/s^2).
\choiceTF
{\True Giá trị cần tìm là $2\,\text{m}$ /s^2.}
{Giá trị cần tìm là $3\,\text{m}$ /s^2.}
{\True Phải xác định đầy đủ các lực hoặc đại lượng liên quan trước khi tính.}
{Có thể bỏ qua điều kiện áp dụng của công thức.}
\loigiai{
\begin{itemchoice}
\item (Đúng) Giá trị cần tìm là $2\,\text{m}$ /s^2. (Vì): kết quả $2\,\text{m}$ /s^2. b) Sai. c) Đúng. d) Sai vì phải kiểm tra điều kiện áp dụng
\item (Sai) Giá trị cần tìm là $3\,\text{m}$ /s^2.
\item (Đúng) Phải xác định đầy đủ các lực hoặc đại lượng liên quan trước khi tính.
\item (Sai) Có thể bỏ qua điều kiện áp dụng của công thức.
\end{itemchoice}
}
\end{ex}

% ===== Câu 15 =====
% ID: AIQ_L10C3_FULL_0339
% Mức: VD
\begin{ex}
% ID: AIQ_L10C3_FULL_0339
% Mức: VD
Vật khối lượng $6\,\text{kg}$ chịu hợp lực $18\,\text{N}$. Tính gia tốc (kết quả theo đơn vị m/s^2).
\choiceTF
{\True Giá trị cần tìm là $3\,\text{m}$ /s^2.}
{Giá trị cần tìm là $4\,\text{m}$ /s^2.}
{\True Phải xác định đầy đủ các lực hoặc đại lượng liên quan trước khi tính.}
{Có thể bỏ qua điều kiện áp dụng của công thức.}
\loigiai{
\begin{itemchoice}
\item (Đúng) Giá trị cần tìm là $3\,\text{m}$ /s^2. (Vì): kết quả $3\,\text{m}$ /s^2. b) Sai. c) Đúng. d) Sai vì phải kiểm tra điều kiện áp dụng
\item (Sai) Giá trị cần tìm là $4\,\text{m}$ /s^2.
\item (Đúng) Phải xác định đầy đủ các lực hoặc đại lượng liên quan trước khi tính.
\item (Sai) Có thể bỏ qua điều kiện áp dụng của công thức.
\end{itemchoice}
}
\end{ex}

% ===== Câu 16 =====
% ID: AIQ_L10C3_FULL_0340
% Mức: TH
\begin{ex}
% ID: AIQ_L10C3_FULL_0340
% Mức: TH
Vật khối lượng $2\,\text{kg}$ chịu hợp lực $8\,\text{N}$. Tính gia tốc (kết quả theo đơn vị m/s^2).
\shortans{4}
\loigiai{
Thay số vào hệ thức của bài học thu được $4\,\text{m}$ /s^2.
}
\end{ex}

% ===== Câu 17 =====
% ID: AIQ_L10C3_FULL_0341
% Mức: TH
\begin{ex}
% ID: AIQ_L10C3_FULL_0341
% Mức: TH
Vật khối lượng $3\,\text{kg}$ chịu hợp lực $3\,\text{N}$. Tính gia tốc (kết quả theo đơn vị m/s^2).
\shortans{1}
\loigiai{
Thay số vào hệ thức của bài học thu được $1\,\text{m}$ /s^2.
}
\end{ex}

% ===== Câu 18 =====
% ID: AIQ_L10C3_FULL_0342
% Mức: TH
\begin{ex}
% ID: AIQ_L10C3_FULL_0342
% Mức: TH
Vật khối lượng $4\,\text{kg}$ chịu hợp lực $8\,\text{N}$. Tính gia tốc (kết quả theo đơn vị m/s^2).
\shortans{2}
\loigiai{
Thay số vào hệ thức của bài học thu được $2\,\text{m}$ /s^2.
}
\end{ex}

% ===== Câu 19 =====
% ID: AIQ_L10C3_FULL_0343
% Mức: VD
\begin{ex}
% ID: AIQ_L10C3_FULL_0343
% Mức: VD
Vật khối lượng $5\,\text{kg}$ chịu hợp lực $15\,\text{N}$. Tính gia tốc (kết quả theo đơn vị m/s^2).
\shortans{3}
\loigiai{
Thay số vào hệ thức của bài học thu được $3\,\text{m}$ /s^2.
}
\end{ex}

% ===== Câu 20 =====
% ID: AIQ_L10C3_FULL_0344
% Mức: VD
\begin{ex}
% ID: AIQ_L10C3_FULL_0344
% Mức: VD
Vật khối lượng $6\,\text{kg}$ chịu hợp lực $24\,\text{N}$. Tính gia tốc (kết quả theo đơn vị m/s^2).
\shortans{4}
\loigiai{
Thay số vào hệ thức của bài học thu được $4\,\text{m}$ /s^2.
}
\end{ex}

% ===== Câu 21 =====
% ID: AIQ_L10C3_FULL_0345
% Mức: VD
\begin{ex}
% ID: AIQ_L10C3_FULL_0345
% Mức: VD
Vật khối lượng $2\,\text{kg}$ chịu hợp lực $2\,\text{N}$. Tính gia tốc (kết quả theo đơn vị m/s^2).
\shortans{1}
\loigiai{
Thay số vào hệ thức của bài học thu được $1\,\text{m}$ /s^2.
}
\end{ex}

% ===== Câu 22 =====
% ID: AIQ_L10C3_FULL_0346
% Mức: VD
\begin{bt}
% ID: AIQ_L10C3_FULL_0346
% Mức: VD
Vật khối lượng $3\,\text{kg}$ chịu hợp lực $6\,\text{N}$. Tính gia tốc (kết quả theo đơn vị m/s^2).
\loigiai{
Xác định các đại lượng đã cho, chọn hệ thức phù hợp và thay số. Kết quả: $2\,\text{m}$ /s^2.
}
\end{bt}

% ===== Câu 23 =====
% ID: AIQ_L10C3_FULL_0347
% Mức: VD
\begin{bt}
% ID: AIQ_L10C3_FULL_0347
% Mức: VD
Vật khối lượng $4\,\text{kg}$ chịu hợp lực $12\,\text{N}$. Tính gia tốc (kết quả theo đơn vị m/s^2).
\loigiai{
Xác định các đại lượng đã cho, chọn hệ thức phù hợp và thay số. Kết quả: $3\,\text{m}$ /s^2.
}
\end{bt}

% ===== Câu 24 =====
% ID: AIQ_L10C3_FULL_0348
% Mức: VD
\begin{bt}
% ID: AIQ_L10C3_FULL_0348
% Mức: VD
Vật khối lượng $5\,\text{kg}$ chịu hợp lực $20\,\text{N}$. Tính gia tốc (kết quả theo đơn vị m/s^2).
\loigiai{
Xác định các đại lượng đã cho, chọn hệ thức phù hợp và thay số. Kết quả: $4\,\text{m}$ /s^2.
}
\end{bt}

% ===== Câu 25 =====
% ID: AIQ_L10C3_FULL_0349
% Mức: VDC
\begin{bt}
% ID: AIQ_L10C3_FULL_0349
% Mức: VDC
Vật khối lượng $6\,\text{kg}$ chịu hợp lực $6\,\text{N}$. Tính gia tốc (kết quả theo đơn vị m/s^2).
\loigiai{
Xác định các đại lượng đã cho, chọn hệ thức phù hợp và thay số. Kết quả: $1\,\text{m}$ /s^2.
}
\end{bt}

% ===== Câu 26 =====
% ID: AIQ_L10C3_FULL_0350
% Mức: VDC
\begin{bt}
% ID: AIQ_L10C3_FULL_0350
% Mức: VDC
Vật khối lượng $2\,\text{kg}$ chịu hợp lực $4\,\text{N}$. Tính gia tốc (kết quả theo đơn vị m/s^2).
\loigiai{
Xác định các đại lượng đã cho, chọn hệ thức phù hợp và thay số. Kết quả: $2\,\text{m}$ /s^2.
}
\end{bt}

% ===== Câu 27 =====
% ID: AIQ_L10C3_FULL_0351
% Mức: VDC
\begin{bt}
% ID: AIQ_L10C3_FULL_0351
% Mức: VDC
Vật khối lượng $3\,\text{kg}$ chịu hợp lực $9\,\text{N}$. Tính gia tốc (kết quả theo đơn vị m/s^2).
\loigiai{
Xác định các đại lượng đã cho, chọn hệ thức phù hợp và thay số. Kết quả: $3\,\text{m}$ /s^2.
}
\end{bt}

% ===== Câu 28 =====
% ID: AIQ_L10C3_FULL_0352
% Mức: NB
\dangbt{Tính gia tốc, lực hoặc khối lượng từ F = ma}
\begin{ex}
% ID: AIQ_L10C3_FULL_0352
% Mức: NB
Vật khối lượng $4\,\text{kg}$ chịu hợp lực $16\,\text{N}$. Tính gia tốc (kết quả theo đơn vị m/s^2).
\choice
{\True 4}
{5}
{3}
{8}
\loigiai{
Áp dụng hệ thức phù hợp của bài học, thay số được $4\,\text{m}$ /s^2.
}
\end{ex}

% ===== Câu 29 =====
% ID: AIQ_L10C3_FULL_0353
% Mức: NB
\begin{ex}
% ID: AIQ_L10C3_FULL_0353
% Mức: NB
Vật khối lượng $5\,\text{kg}$ chịu hợp lực $5\,\text{N}$. Tính gia tốc (kết quả theo đơn vị m/s^2).
\choice
{2}
{\True 1}
{0}
{2}
\loigiai{
Áp dụng hệ thức phù hợp của bài học, thay số được $1\,\text{m}$ /s^2.
}
\end{ex}

% ===== Câu 30 =====
% ID: AIQ_L10C3_FULL_0354
% Mức: NB
\begin{ex}
% ID: AIQ_L10C3_FULL_0354
% Mức: NB
Vật khối lượng $6\,\text{kg}$ chịu hợp lực $12\,\text{N}$. Tính gia tốc (kết quả theo đơn vị m/s^2).
\choice
{1}
{3}
{\True 2}
{4}
\loigiai{
Áp dụng hệ thức phù hợp của bài học, thay số được $2\,\text{m}$ /s^2.
}
\end{ex}

% ===== Câu 31 =====
% ID: AIQ_L10C3_FULL_0355
% Mức: NB
\begin{ex}
% ID: AIQ_L10C3_FULL_0355
% Mức: NB
Vật khối lượng $2\,\text{kg}$ chịu hợp lực $6\,\text{N}$. Tính gia tốc (kết quả theo đơn vị m/s^2).
\choice
{6}
{4}
{2}
{\True 3}
\loigiai{
Áp dụng hệ thức phù hợp của bài học, thay số được $3\,\text{m}$ /s^2.
}
\end{ex}

% ===== Câu 32 =====
% ID: AIQ_L10C3_FULL_0356
% Mức: TH
\begin{ex}
% ID: AIQ_L10C3_FULL_0356
% Mức: TH
Vật khối lượng $3\,\text{kg}$ chịu hợp lực $12\,\text{N}$. Tính gia tốc (kết quả theo đơn vị m/s^2).
\choice
{\True 4}
{5}
{3}
{8}
\loigiai{
Áp dụng hệ thức phù hợp của bài học, thay số được $4\,\text{m}$ /s^2.
}
\end{ex}

% ===== Câu 33 =====
% ID: AIQ_L10C3_FULL_0357
% Mức: TH
\begin{ex}
% ID: AIQ_L10C3_FULL_0357
% Mức: TH
Vật khối lượng $4\,\text{kg}$ chịu hợp lực $4\,\text{N}$. Tính gia tốc (kết quả theo đơn vị m/s^2).
\choice
{2}
{\True 1}
{0}
{2}
\loigiai{
Áp dụng hệ thức phù hợp của bài học, thay số được $1\,\text{m}$ /s^2.
}
\end{ex}

% ===== Câu 34 =====
% ID: AIQ_L10C3_FULL_0358
% Mức: TH
\begin{ex}
% ID: AIQ_L10C3_FULL_0358
% Mức: TH
Vật khối lượng $5\,\text{kg}$ chịu hợp lực $10\,\text{N}$. Tính gia tốc (kết quả theo đơn vị m/s^2).
\choice
{1}
{3}
{\True 2}
{4}
\loigiai{
Áp dụng hệ thức phù hợp của bài học, thay số được $2\,\text{m}$ /s^2.
}
\end{ex}

% ===== Câu 35 =====
% ID: AIQ_L10C3_FULL_0359
% Mức: TH
\begin{ex}
% ID: AIQ_L10C3_FULL_0359
% Mức: TH
Vật khối lượng $6\,\text{kg}$ chịu hợp lực $18\,\text{N}$. Tính gia tốc (kết quả theo đơn vị m/s^2).
\choice
{6}
{4}
{2}
{\True 3}
\loigiai{
Áp dụng hệ thức phù hợp của bài học, thay số được $3\,\text{m}$ /s^2.
}
\end{ex}

% ===== Câu 36 =====
% ID: AIQ_L10C3_FULL_0360
% Mức: VD
\begin{ex}
% ID: AIQ_L10C3_FULL_0360
% Mức: VD
Vật khối lượng $2\,\text{kg}$ chịu hợp lực $8\,\text{N}$. Tính gia tốc (kết quả theo đơn vị m/s^2).
\choice
{\True 4}
{5}
{3}
{8}
\loigiai{
Áp dụng hệ thức phù hợp của bài học, thay số được $4\,\text{m}$ /s^2.
}
\end{ex}

% ===== Câu 37 =====
% ID: AIQ_L10C3_FULL_0361
% Mức: VD
\begin{ex}
% ID: AIQ_L10C3_FULL_0361
% Mức: VD
Vật khối lượng $3\,\text{kg}$ chịu hợp lực $3\,\text{N}$. Tính gia tốc (kết quả theo đơn vị m/s^2).
\choice
{2}
{\True 1}
{0}
{2}
\loigiai{
Áp dụng hệ thức phù hợp của bài học, thay số được $1\,\text{m}$ /s^2.
}
\end{ex}

% ===== Câu 38 =====
% ID: AIQ_L10C3_FULL_0362
% Mức: TH
\begin{ex}
% ID: AIQ_L10C3_FULL_0362
% Mức: TH
Vật khối lượng $4\,\text{kg}$ chịu hợp lực $8\,\text{N}$. Tính gia tốc (kết quả theo đơn vị m/s^2).
\choiceTF
{\True Giá trị cần tìm là $2\,\text{m}$ /s^2.}
{Giá trị cần tìm là $3\,\text{m}$ /s^2.}
{\True Phải xác định đầy đủ các lực hoặc đại lượng liên quan trước khi tính.}
{Có thể bỏ qua điều kiện áp dụng của công thức.}
\loigiai{
\begin{itemchoice}
\item (Đúng) Giá trị cần tìm là $2\,\text{m}$ /s^2. (Vì): kết quả $2\,\text{m}$ /s^2. b) Sai. c) Đúng. d) Sai vì phải kiểm tra điều kiện áp dụng
\item (Sai) Giá trị cần tìm là $3\,\text{m}$ /s^2.
\item (Đúng) Phải xác định đầy đủ các lực hoặc đại lượng liên quan trước khi tính.
\item (Sai) Có thể bỏ qua điều kiện áp dụng của công thức.
\end{itemchoice}
}
\end{ex}

% ===== Câu 39 =====
% ID: AIQ_L10C3_FULL_0363
% Mức: TH
\begin{ex}
% ID: AIQ_L10C3_FULL_0363
% Mức: TH
Vật khối lượng $5\,\text{kg}$ chịu hợp lực $15\,\text{N}$. Tính gia tốc (kết quả theo đơn vị m/s^2).
\choiceTF
{\True Giá trị cần tìm là $3\,\text{m}$ /s^2.}
{Giá trị cần tìm là $4\,\text{m}$ /s^2.}
{\True Phải xác định đầy đủ các lực hoặc đại lượng liên quan trước khi tính.}
{Có thể bỏ qua điều kiện áp dụng của công thức.}
\loigiai{
\begin{itemchoice}
\item (Đúng) Giá trị cần tìm là $3\,\text{m}$ /s^2. (Vì): kết quả $3\,\text{m}$ /s^2. b) Sai. c) Đúng. d) Sai vì phải kiểm tra điều kiện áp dụng
\item (Sai) Giá trị cần tìm là $4\,\text{m}$ /s^2.
\item (Đúng) Phải xác định đầy đủ các lực hoặc đại lượng liên quan trước khi tính.
\item (Sai) Có thể bỏ qua điều kiện áp dụng của công thức.
\end{itemchoice}
}
\end{ex}

% ===== Câu 40 =====
% ID: AIQ_L10C3_FULL_0364
% Mức: VD
\begin{ex}
% ID: AIQ_L10C3_FULL_0364
% Mức: VD
Vật khối lượng $6\,\text{kg}$ chịu hợp lực $24\,\text{N}$. Tính gia tốc (kết quả theo đơn vị m/s^2).
\choiceTF
{\True Giá trị cần tìm là $4\,\text{m}$ /s^2.}
{Giá trị cần tìm là $5\,\text{m}$ /s^2.}
{\True Phải xác định đầy đủ các lực hoặc đại lượng liên quan trước khi tính.}
{Có thể bỏ qua điều kiện áp dụng của công thức.}
\loigiai{
\begin{itemchoice}
\item (Đúng) Giá trị cần tìm là $4\,\text{m}$ /s^2. (Vì): kết quả $4\,\text{m}$ /s^2. b) Sai. c) Đúng. d) Sai vì phải kiểm tra điều kiện áp dụng
\item (Sai) Giá trị cần tìm là $5\,\text{m}$ /s^2.
\item (Đúng) Phải xác định đầy đủ các lực hoặc đại lượng liên quan trước khi tính.
\item (Sai) Có thể bỏ qua điều kiện áp dụng của công thức.
\end{itemchoice}
}
\end{ex}

% ===== Câu 41 =====
% ID: AIQ_L10C3_FULL_0365
% Mức: VD
\begin{ex}
% ID: AIQ_L10C3_FULL_0365
% Mức: VD
Vật khối lượng $2\,\text{kg}$ chịu hợp lực $2\,\text{N}$. Tính gia tốc (kết quả theo đơn vị m/s^2).
\choiceTF
{\True Giá trị cần tìm là $1\,\text{m}$ /s^2.}
{Giá trị cần tìm là $2\,\text{m}$ /s^2.}
{\True Phải xác định đầy đủ các lực hoặc đại lượng liên quan trước khi tính.}
{Có thể bỏ qua điều kiện áp dụng của công thức.}
\loigiai{
\begin{itemchoice}
\item (Đúng) Giá trị cần tìm là $1\,\text{m}$ /s^2. (Vì): kết quả $1\,\text{m}$ /s^2. b) Sai. c) Đúng. d) Sai vì phải kiểm tra điều kiện áp dụng
\item (Sai) Giá trị cần tìm là $2\,\text{m}$ /s^2.
\item (Đúng) Phải xác định đầy đủ các lực hoặc đại lượng liên quan trước khi tính.
\item (Sai) Có thể bỏ qua điều kiện áp dụng của công thức.
\end{itemchoice}
}
\end{ex}

% ===== Câu 42 =====
% ID: AIQ_L10C3_FULL_0366
% Mức: VD
\begin{ex}
% ID: AIQ_L10C3_FULL_0366
% Mức: VD
Vật khối lượng $3\,\text{kg}$ chịu hợp lực $6\,\text{N}$. Tính gia tốc (kết quả theo đơn vị m/s^2).
\choiceTF
{\True Giá trị cần tìm là $2\,\text{m}$ /s^2.}
{Giá trị cần tìm là $3\,\text{m}$ /s^2.}
{\True Phải xác định đầy đủ các lực hoặc đại lượng liên quan trước khi tính.}
{Có thể bỏ qua điều kiện áp dụng của công thức.}
\loigiai{
\begin{itemchoice}
\item (Đúng) Giá trị cần tìm là $2\,\text{m}$ /s^2. (Vì): kết quả $2\,\text{m}$ /s^2. b) Sai. c) Đúng. d) Sai vì phải kiểm tra điều kiện áp dụng
\item (Sai) Giá trị cần tìm là $3\,\text{m}$ /s^2.
\item (Đúng) Phải xác định đầy đủ các lực hoặc đại lượng liên quan trước khi tính.
\item (Sai) Có thể bỏ qua điều kiện áp dụng của công thức.
\end{itemchoice}
}
\end{ex}

% ===== Câu 43 =====
% ID: AIQ_L10C3_FULL_0367
% Mức: TH
\begin{ex}
% ID: AIQ_L10C3_FULL_0367
% Mức: TH
Vật khối lượng $4\,\text{kg}$ chịu hợp lực $12\,\text{N}$. Tính gia tốc (kết quả theo đơn vị m/s^2).
\shortans{3}
\loigiai{
Thay số vào hệ thức của bài học thu được $3\,\text{m}$ /s^2.
}
\end{ex}

% ===== Câu 44 =====
% ID: AIQ_L10C3_FULL_0368
% Mức: TH
\begin{ex}
% ID: AIQ_L10C3_FULL_0368
% Mức: TH
Vật khối lượng $5\,\text{kg}$ chịu hợp lực $20\,\text{N}$. Tính gia tốc (kết quả theo đơn vị m/s^2).
\shortans{4}
\loigiai{
Thay số vào hệ thức của bài học thu được $4\,\text{m}$ /s^2.
}
\end{ex}

% ===== Câu 45 =====
% ID: AIQ_L10C3_FULL_0369
% Mức: TH
\begin{ex}
% ID: AIQ_L10C3_FULL_0369
% Mức: TH
Vật khối lượng $6\,\text{kg}$ chịu hợp lực $6\,\text{N}$. Tính gia tốc (kết quả theo đơn vị m/s^2).
\shortans{1}
\loigiai{
Thay số vào hệ thức của bài học thu được $1\,\text{m}$ /s^2.
}
\end{ex}

% ===== Câu 46 =====
% ID: AIQ_L10C3_FULL_0370
% Mức: VD
\begin{ex}
% ID: AIQ_L10C3_FULL_0370
% Mức: VD
Vật khối lượng $2\,\text{kg}$ chịu hợp lực $4\,\text{N}$. Tính gia tốc (kết quả theo đơn vị m/s^2).
\shortans{2}
\loigiai{
Thay số vào hệ thức của bài học thu được $2\,\text{m}$ /s^2.
}
\end{ex}

% ===== Câu 47 =====
% ID: AIQ_L10C3_FULL_0371
% Mức: VD
\begin{ex}
% ID: AIQ_L10C3_FULL_0371
% Mức: VD
Vật khối lượng $3\,\text{kg}$ chịu hợp lực $9\,\text{N}$. Tính gia tốc (kết quả theo đơn vị m/s^2).
\shortans{3}
\loigiai{
Thay số vào hệ thức của bài học thu được $3\,\text{m}$ /s^2.
}
\end{ex}

% ===== Câu 48 =====
% ID: AIQ_L10C3_FULL_0372
% Mức: VD
\begin{ex}
% ID: AIQ_L10C3_FULL_0372
% Mức: VD
Vật khối lượng $4\,\text{kg}$ chịu hợp lực $16\,\text{N}$. Tính gia tốc (kết quả theo đơn vị m/s^2).
\shortans{4}
\loigiai{
Thay số vào hệ thức của bài học thu được $4\,\text{m}$ /s^2.
}
\end{ex}

% ===== Câu 49 =====
% ID: AIQ_L10C3_FULL_0373
% Mức: VD
\begin{bt}
% ID: AIQ_L10C3_FULL_0373
% Mức: VD
Vật khối lượng $5\,\text{kg}$ chịu hợp lực $5\,\text{N}$. Tính gia tốc (kết quả theo đơn vị m/s^2).
\loigiai{
Xác định các đại lượng đã cho, chọn hệ thức phù hợp và thay số. Kết quả: $1\,\text{m}$ /s^2.
}
\end{bt}

% ===== Câu 50 =====
% ID: AIQ_L10C3_FULL_0374
% Mức: VD
\begin{bt}
% ID: AIQ_L10C3_FULL_0374
% Mức: VD
Vật khối lượng $6\,\text{kg}$ chịu hợp lực $12\,\text{N}$. Tính gia tốc (kết quả theo đơn vị m/s^2).
\loigiai{
Xác định các đại lượng đã cho, chọn hệ thức phù hợp và thay số. Kết quả: $2\,\text{m}$ /s^2.
}
\end{bt}

% ===== Câu 51 =====
% ID: AIQ_L10C3_FULL_0375
% Mức: VD
\begin{bt}
% ID: AIQ_L10C3_FULL_0375
% Mức: VD
Vật khối lượng $2\,\text{kg}$ chịu hợp lực $6\,\text{N}$. Tính gia tốc (kết quả theo đơn vị m/s^2).
\loigiai{
Xác định các đại lượng đã cho, chọn hệ thức phù hợp và thay số. Kết quả: $3\,\text{m}$ /s^2.
}
\end{bt}

% ===== Câu 52 =====
% ID: AIQ_L10C3_FULL_0376
% Mức: VDC
\begin{bt}
% ID: AIQ_L10C3_FULL_0376
% Mức: VDC
Vật khối lượng $3\,\text{kg}$ chịu hợp lực $12\,\text{N}$. Tính gia tốc (kết quả theo đơn vị m/s^2).
\loigiai{
Xác định các đại lượng đã cho, chọn hệ thức phù hợp và thay số. Kết quả: $4\,\text{m}$ /s^2.
}
\end{bt}

% ===== Câu 53 =====
% ID: AIQ_L10C3_FULL_0377
% Mức: VDC
\begin{bt}
% ID: AIQ_L10C3_FULL_0377
% Mức: VDC
Vật khối lượng $4\,\text{kg}$ chịu hợp lực $4\,\text{N}$. Tính gia tốc (kết quả theo đơn vị m/s^2).
\loigiai{
Xác định các đại lượng đã cho, chọn hệ thức phù hợp và thay số. Kết quả: $1\,\text{m}$ /s^2.
}
\end{bt}

% ===== Câu 54 =====
% ID: AIQ_L10C3_FULL_0378
% Mức: VDC
\begin{bt}
% ID: AIQ_L10C3_FULL_0378
% Mức: VDC
Vật khối lượng $5\,\text{kg}$ chịu hợp lực $10\,\text{N}$. Tính gia tốc (kết quả theo đơn vị m/s^2).
\loigiai{
Xác định các đại lượng đã cho, chọn hệ thức phù hợp và thay số. Kết quả: $2\,\text{m}$ /s^2.
}
\end{bt}

% ===== Câu 55 =====
% ID: AIQ_L10C3_FULL_0379
% Mức: NB
\dangbt{Vận dụng định luật II Newton cho chuyển động ngang}
\begin{ex}
% ID: AIQ_L10C3_FULL_0379
% Mức: NB
Vật khối lượng $6\,\text{kg}$ chịu hợp lực $18\,\text{N}$. Tính gia tốc (kết quả theo đơn vị m/s^2).
\choice
{\True 3}
{4}
{2}
{6}
\loigiai{
Áp dụng hệ thức phù hợp của bài học, thay số được $3\,\text{m}$ /s^2.
}
\end{ex}

% ===== Câu 56 =====
% ID: AIQ_L10C3_FULL_0380
% Mức: NB
\begin{ex}
% ID: AIQ_L10C3_FULL_0380
% Mức: NB
Vật khối lượng $2\,\text{kg}$ chịu hợp lực $8\,\text{N}$. Tính gia tốc (kết quả theo đơn vị m/s^2).
\choice
{5}
{\True 4}
{3}
{8}
\loigiai{
Áp dụng hệ thức phù hợp của bài học, thay số được $4\,\text{m}$ /s^2.
}
\end{ex}

% ===== Câu 57 =====
% ID: AIQ_L10C3_FULL_0381
% Mức: NB
\begin{ex}
% ID: AIQ_L10C3_FULL_0381
% Mức: NB
Vật khối lượng $3\,\text{kg}$ chịu hợp lực $3\,\text{N}$. Tính gia tốc (kết quả theo đơn vị m/s^2).
\choice
{0}
{2}
{\True 1}
{2}
\loigiai{
Áp dụng hệ thức phù hợp của bài học, thay số được $1\,\text{m}$ /s^2.
}
\end{ex}

% ===== Câu 58 =====
% ID: AIQ_L10C3_FULL_0382
% Mức: NB
\begin{ex}
% ID: AIQ_L10C3_FULL_0382
% Mức: NB
Vật khối lượng $4\,\text{kg}$ chịu hợp lực $8\,\text{N}$. Tính gia tốc (kết quả theo đơn vị m/s^2).
\choice
{4}
{3}
{1}
{\True 2}
\loigiai{
Áp dụng hệ thức phù hợp của bài học, thay số được $2\,\text{m}$ /s^2.
}
\end{ex}

% ===== Câu 59 =====
% ID: AIQ_L10C3_FULL_0383
% Mức: TH
\begin{ex}
% ID: AIQ_L10C3_FULL_0383
% Mức: TH
Vật khối lượng $5\,\text{kg}$ chịu hợp lực $15\,\text{N}$. Tính gia tốc (kết quả theo đơn vị m/s^2).
\choice
{\True 3}
{4}
{2}
{6}
\loigiai{
Áp dụng hệ thức phù hợp của bài học, thay số được $3\,\text{m}$ /s^2.
}
\end{ex}

% ===== Câu 60 =====
% ID: AIQ_L10C3_FULL_0384
% Mức: TH
\begin{ex}
% ID: AIQ_L10C3_FULL_0384
% Mức: TH
Vật khối lượng $6\,\text{kg}$ chịu hợp lực $24\,\text{N}$. Tính gia tốc (kết quả theo đơn vị m/s^2).
\choice
{5}
{\True 4}
{3}
{8}
\loigiai{
Áp dụng hệ thức phù hợp của bài học, thay số được $4\,\text{m}$ /s^2.
}
\end{ex}

% ===== Câu 61 =====
% ID: AIQ_L10C3_FULL_0385
% Mức: TH
\begin{ex}
% ID: AIQ_L10C3_FULL_0385
% Mức: TH
Vật khối lượng $2\,\text{kg}$ chịu hợp lực $2\,\text{N}$. Tính gia tốc (kết quả theo đơn vị m/s^2).
\choice
{0}
{2}
{\True 1}
{2}
\loigiai{
Áp dụng hệ thức phù hợp của bài học, thay số được $1\,\text{m}$ /s^2.
}
\end{ex}

% ===== Câu 62 =====
% ID: AIQ_L10C3_FULL_0386
% Mức: TH
\begin{ex}
% ID: AIQ_L10C3_FULL_0386
% Mức: TH
Vật khối lượng $3\,\text{kg}$ chịu hợp lực $6\,\text{N}$. Tính gia tốc (kết quả theo đơn vị m/s^2).
\choice
{4}
{3}
{1}
{\True 2}
\loigiai{
Áp dụng hệ thức phù hợp của bài học, thay số được $2\,\text{m}$ /s^2.
}
\end{ex}

% ===== Câu 63 =====
% ID: AIQ_L10C3_FULL_0387
% Mức: VD
\begin{ex}
% ID: AIQ_L10C3_FULL_0387
% Mức: VD
Vật khối lượng $4\,\text{kg}$ chịu hợp lực $12\,\text{N}$. Tính gia tốc (kết quả theo đơn vị m/s^2).
\choice
{\True 3}
{4}
{2}
{6}
\loigiai{
Áp dụng hệ thức phù hợp của bài học, thay số được $3\,\text{m}$ /s^2.
}
\end{ex}

% ===== Câu 64 =====
% ID: AIQ_L10C3_FULL_0388
% Mức: VD
\begin{ex}
% ID: AIQ_L10C3_FULL_0388
% Mức: VD
Vật khối lượng $5\,\text{kg}$ chịu hợp lực $20\,\text{N}$. Tính gia tốc (kết quả theo đơn vị m/s^2).
\choice
{5}
{\True 4}
{3}
{8}
\loigiai{
Áp dụng hệ thức phù hợp của bài học, thay số được $4\,\text{m}$ /s^2.
}
\end{ex}

% ===== Câu 65 =====
% ID: AIQ_L10C3_FULL_0389
% Mức: TH
\begin{ex}
% ID: AIQ_L10C3_FULL_0389
% Mức: TH
Vật khối lượng $6\,\text{kg}$ chịu hợp lực $6\,\text{N}$. Tính gia tốc (kết quả theo đơn vị m/s^2).
\choiceTF
{\True Giá trị cần tìm là $1\,\text{m}$ /s^2.}
{Giá trị cần tìm là $2\,\text{m}$ /s^2.}
{\True Phải xác định đầy đủ các lực hoặc đại lượng liên quan trước khi tính.}
{Có thể bỏ qua điều kiện áp dụng của công thức.}
\loigiai{
\begin{itemchoice}
\item (Đúng) Giá trị cần tìm là $1\,\text{m}$ /s^2. (Vì): kết quả $1\,\text{m}$ /s^2. b) Sai. c) Đúng. d) Sai vì phải kiểm tra điều kiện áp dụng
\item (Sai) Giá trị cần tìm là $2\,\text{m}$ /s^2.
\item (Đúng) Phải xác định đầy đủ các lực hoặc đại lượng liên quan trước khi tính.
\item (Sai) Có thể bỏ qua điều kiện áp dụng của công thức.
\end{itemchoice}
}
\end{ex}

% ===== Câu 66 =====
% ID: AIQ_L10C3_FULL_0390
% Mức: TH
\begin{ex}
% ID: AIQ_L10C3_FULL_0390
% Mức: TH
Vật khối lượng $2\,\text{kg}$ chịu hợp lực $4\,\text{N}$. Tính gia tốc (kết quả theo đơn vị m/s^2).
\choiceTF
{\True Giá trị cần tìm là $2\,\text{m}$ /s^2.}
{Giá trị cần tìm là $3\,\text{m}$ /s^2.}
{\True Phải xác định đầy đủ các lực hoặc đại lượng liên quan trước khi tính.}
{Có thể bỏ qua điều kiện áp dụng của công thức.}
\loigiai{
\begin{itemchoice}
\item (Đúng) Giá trị cần tìm là $2\,\text{m}$ /s^2. (Vì): kết quả $2\,\text{m}$ /s^2. b) Sai. c) Đúng. d) Sai vì phải kiểm tra điều kiện áp dụng
\item (Sai) Giá trị cần tìm là $3\,\text{m}$ /s^2.
\item (Đúng) Phải xác định đầy đủ các lực hoặc đại lượng liên quan trước khi tính.
\item (Sai) Có thể bỏ qua điều kiện áp dụng của công thức.
\end{itemchoice}
}
\end{ex}

% ===== Câu 67 =====
% ID: AIQ_L10C3_FULL_0391
% Mức: VD
\begin{ex}
% ID: AIQ_L10C3_FULL_0391
% Mức: VD
Vật khối lượng $3\,\text{kg}$ chịu hợp lực $9\,\text{N}$. Tính gia tốc (kết quả theo đơn vị m/s^2).
\choiceTF
{\True Giá trị cần tìm là $3\,\text{m}$ /s^2.}
{Giá trị cần tìm là $4\,\text{m}$ /s^2.}
{\True Phải xác định đầy đủ các lực hoặc đại lượng liên quan trước khi tính.}
{Có thể bỏ qua điều kiện áp dụng của công thức.}
\loigiai{
\begin{itemchoice}
\item (Đúng) Giá trị cần tìm là $3\,\text{m}$ /s^2. (Vì): kết quả $3\,\text{m}$ /s^2. b) Sai. c) Đúng. d) Sai vì phải kiểm tra điều kiện áp dụng
\item (Sai) Giá trị cần tìm là $4\,\text{m}$ /s^2.
\item (Đúng) Phải xác định đầy đủ các lực hoặc đại lượng liên quan trước khi tính.
\item (Sai) Có thể bỏ qua điều kiện áp dụng của công thức.
\end{itemchoice}
}
\end{ex}

% ===== Câu 68 =====
% ID: AIQ_L10C3_FULL_0392
% Mức: VD
\begin{ex}
% ID: AIQ_L10C3_FULL_0392
% Mức: VD
Vật khối lượng $4\,\text{kg}$ chịu hợp lực $16\,\text{N}$. Tính gia tốc (kết quả theo đơn vị m/s^2).
\choiceTF
{\True Giá trị cần tìm là $4\,\text{m}$ /s^2.}
{Giá trị cần tìm là $5\,\text{m}$ /s^2.}
{\True Phải xác định đầy đủ các lực hoặc đại lượng liên quan trước khi tính.}
{Có thể bỏ qua điều kiện áp dụng của công thức.}
\loigiai{
\begin{itemchoice}
\item (Đúng) Giá trị cần tìm là $4\,\text{m}$ /s^2. (Vì): kết quả $4\,\text{m}$ /s^2. b) Sai. c) Đúng. d) Sai vì phải kiểm tra điều kiện áp dụng
\item (Sai) Giá trị cần tìm là $5\,\text{m}$ /s^2.
\item (Đúng) Phải xác định đầy đủ các lực hoặc đại lượng liên quan trước khi tính.
\item (Sai) Có thể bỏ qua điều kiện áp dụng của công thức.
\end{itemchoice}
}
\end{ex}

% ===== Câu 69 =====
% ID: AIQ_L10C3_FULL_0393
% Mức: VD
\begin{ex}
% ID: AIQ_L10C3_FULL_0393
% Mức: VD
Vật khối lượng $5\,\text{kg}$ chịu hợp lực $5\,\text{N}$. Tính gia tốc (kết quả theo đơn vị m/s^2).
\choiceTF
{\True Giá trị cần tìm là $1\,\text{m}$ /s^2.}
{Giá trị cần tìm là $2\,\text{m}$ /s^2.}
{\True Phải xác định đầy đủ các lực hoặc đại lượng liên quan trước khi tính.}
{Có thể bỏ qua điều kiện áp dụng của công thức.}
\loigiai{
\begin{itemchoice}
\item (Đúng) Giá trị cần tìm là $1\,\text{m}$ /s^2. (Vì): kết quả $1\,\text{m}$ /s^2. b) Sai. c) Đúng. d) Sai vì phải kiểm tra điều kiện áp dụng
\item (Sai) Giá trị cần tìm là $2\,\text{m}$ /s^2.
\item (Đúng) Phải xác định đầy đủ các lực hoặc đại lượng liên quan trước khi tính.
\item (Sai) Có thể bỏ qua điều kiện áp dụng của công thức.
\end{itemchoice}
}
\end{ex}

% ===== Câu 70 =====
% ID: AIQ_L10C3_FULL_0394
% Mức: TH
\begin{ex}
% ID: AIQ_L10C3_FULL_0394
% Mức: TH
Vật khối lượng $6\,\text{kg}$ chịu hợp lực $12\,\text{N}$. Tính gia tốc (kết quả theo đơn vị m/s^2).
\shortans{2}
\loigiai{
Thay số vào hệ thức của bài học thu được $2\,\text{m}$ /s^2.
}
\end{ex}

% ===== Câu 71 =====
% ID: AIQ_L10C3_FULL_0395
% Mức: TH
\begin{ex}
% ID: AIQ_L10C3_FULL_0395
% Mức: TH
Vật khối lượng $2\,\text{kg}$ chịu hợp lực $6\,\text{N}$. Tính gia tốc (kết quả theo đơn vị m/s^2).
\shortans{3}
\loigiai{
Thay số vào hệ thức của bài học thu được $3\,\text{m}$ /s^2.
}
\end{ex}

% ===== Câu 72 =====
% ID: AIQ_L10C3_FULL_0396
% Mức: TH
\begin{ex}
% ID: AIQ_L10C3_FULL_0396
% Mức: TH
Vật khối lượng $3\,\text{kg}$ chịu hợp lực $12\,\text{N}$. Tính gia tốc (kết quả theo đơn vị m/s^2).
\shortans{4}
\loigiai{
Thay số vào hệ thức của bài học thu được $4\,\text{m}$ /s^2.
}
\end{ex}

% ===== Câu 73 =====
% ID: AIQ_L10C3_FULL_0397
% Mức: VD
\begin{ex}
% ID: AIQ_L10C3_FULL_0397
% Mức: VD
Vật khối lượng $4\,\text{kg}$ chịu hợp lực $4\,\text{N}$. Tính gia tốc (kết quả theo đơn vị m/s^2).
\shortans{1}
\loigiai{
Thay số vào hệ thức của bài học thu được $1\,\text{m}$ /s^2.
}
\end{ex}

% ===== Câu 74 =====
% ID: AIQ_L10C3_FULL_0398
% Mức: VD
\begin{ex}
% ID: AIQ_L10C3_FULL_0398
% Mức: VD
Vật khối lượng $5\,\text{kg}$ chịu hợp lực $10\,\text{N}$. Tính gia tốc (kết quả theo đơn vị m/s^2).
\shortans{2}
\loigiai{
Thay số vào hệ thức của bài học thu được $2\,\text{m}$ /s^2.
}
\end{ex}

% ===== Câu 75 =====
% ID: AIQ_L10C3_FULL_0399
% Mức: VD
\begin{ex}
% ID: AIQ_L10C3_FULL_0399
% Mức: VD
Vật khối lượng $6\,\text{kg}$ chịu hợp lực $18\,\text{N}$. Tính gia tốc (kết quả theo đơn vị m/s^2).
\shortans{3}
\loigiai{
Thay số vào hệ thức của bài học thu được $3\,\text{m}$ /s^2.
}
\end{ex}

% ===== Câu 76 =====
% ID: AIQ_L10C3_FULL_0400
% Mức: VD
\begin{bt}
% ID: AIQ_L10C3_FULL_0400
% Mức: VD
Vật khối lượng $2\,\text{kg}$ chịu hợp lực $8\,\text{N}$. Tính gia tốc (kết quả theo đơn vị m/s^2).
\loigiai{
Xác định các đại lượng đã cho, chọn hệ thức phù hợp và thay số. Kết quả: $4\,\text{m}$ /s^2.
}
\end{bt}

% ===== Câu 77 =====
% ID: AIQ_L10C3_FULL_0401
% Mức: VD
\begin{bt}
% ID: AIQ_L10C3_FULL_0401
% Mức: VD
Vật khối lượng $3\,\text{kg}$ chịu hợp lực $3\,\text{N}$. Tính gia tốc (kết quả theo đơn vị m/s^2).
\loigiai{
Xác định các đại lượng đã cho, chọn hệ thức phù hợp và thay số. Kết quả: $1\,\text{m}$ /s^2.
}
\end{bt}

% ===== Câu 78 =====
% ID: AIQ_L10C3_FULL_0402
% Mức: VD
\begin{bt}
% ID: AIQ_L10C3_FULL_0402
% Mức: VD
Vật khối lượng $4\,\text{kg}$ chịu hợp lực $8\,\text{N}$. Tính gia tốc (kết quả theo đơn vị m/s^2).
\loigiai{
Xác định các đại lượng đã cho, chọn hệ thức phù hợp và thay số. Kết quả: $2\,\text{m}$ /s^2.
}
\end{bt}

% ===== Câu 79 =====
% ID: AIQ_L10C3_FULL_0403
% Mức: VDC
\begin{bt}
% ID: AIQ_L10C3_FULL_0403
% Mức: VDC
Vật khối lượng $5\,\text{kg}$ chịu hợp lực $15\,\text{N}$. Tính gia tốc (kết quả theo đơn vị m/s^2).
\loigiai{
Xác định các đại lượng đã cho, chọn hệ thức phù hợp và thay số. Kết quả: $3\,\text{m}$ /s^2.
}
\end{bt}

% ===== Câu 80 =====
% ID: AIQ_L10C3_FULL_0404
% Mức: VDC
\begin{bt}
% ID: AIQ_L10C3_FULL_0404
% Mức: VDC
Vật khối lượng $6\,\text{kg}$ chịu hợp lực $24\,\text{N}$. Tính gia tốc (kết quả theo đơn vị m/s^2).
\loigiai{
Xác định các đại lượng đã cho, chọn hệ thức phù hợp và thay số. Kết quả: $4\,\text{m}$ /s^2.
}
\end{bt}

% ===== Câu 81 =====
% ID: AIQ_L10C3_FULL_0405
% Mức: VDC
\begin{bt}
% ID: AIQ_L10C3_FULL_0405
% Mức: VDC
Vật khối lượng $2\,\text{kg}$ chịu hợp lực $2\,\text{N}$. Tính gia tốc (kết quả theo đơn vị m/s^2).
\loigiai{
Xác định các đại lượng đã cho, chọn hệ thức phù hợp và thay số. Kết quả: $1\,\text{m}$ /s^2.
}
\end{bt}

% ===== Câu 82 =====
% ID: AIQ_L10C3_FULL_0406
% Mức: NB
\dangbt{Vận dụng định luật II Newton cho chuyển động trên mặt phẳng nghiêng}
\begin{ex}
% ID: AIQ_L10C3_FULL_0406
% Mức: NB
Vật khối lượng $3\,\text{kg}$ chịu hợp lực $6\,\text{N}$. Tính gia tốc (kết quả theo đơn vị m/s^2).
\choice
{\True 2}
{3}
{1}
{4}
\loigiai{
Áp dụng hệ thức phù hợp của bài học, thay số được $2\,\text{m}$ /s^2.
}
\end{ex}

% ===== Câu 83 =====
% ID: AIQ_L10C3_FULL_0407
% Mức: NB
\begin{ex}
% ID: AIQ_L10C3_FULL_0407
% Mức: NB
Vật khối lượng $4\,\text{kg}$ chịu hợp lực $12\,\text{N}$. Tính gia tốc (kết quả theo đơn vị m/s^2).
\choice
{4}
{\True 3}
{2}
{6}
\loigiai{
Áp dụng hệ thức phù hợp của bài học, thay số được $3\,\text{m}$ /s^2.
}
\end{ex}

% ===== Câu 84 =====
% ID: AIQ_L10C3_FULL_0408
% Mức: NB
\begin{ex}
% ID: AIQ_L10C3_FULL_0408
% Mức: NB
Vật khối lượng $5\,\text{kg}$ chịu hợp lực $20\,\text{N}$. Tính gia tốc (kết quả theo đơn vị m/s^2).
\choice
{3}
{5}
{\True 4}
{8}
\loigiai{
Áp dụng hệ thức phù hợp của bài học, thay số được $4\,\text{m}$ /s^2.
}
\end{ex}

% ===== Câu 85 =====
% ID: AIQ_L10C3_FULL_0409
% Mức: NB
\begin{ex}
% ID: AIQ_L10C3_FULL_0409
% Mức: NB
Vật khối lượng $6\,\text{kg}$ chịu hợp lực $6\,\text{N}$. Tính gia tốc (kết quả theo đơn vị m/s^2).
\choice
{2}
{2}
{0}
{\True 1}
\loigiai{
Áp dụng hệ thức phù hợp của bài học, thay số được $1\,\text{m}$ /s^2.
}
\end{ex}

% ===== Câu 86 =====
% ID: AIQ_L10C3_FULL_0410
% Mức: TH
\begin{ex}
% ID: AIQ_L10C3_FULL_0410
% Mức: TH
Vật khối lượng $2\,\text{kg}$ chịu hợp lực $4\,\text{N}$. Tính gia tốc (kết quả theo đơn vị m/s^2).
\choice
{\True 2}
{3}
{1}
{4}
\loigiai{
Áp dụng hệ thức phù hợp của bài học, thay số được $2\,\text{m}$ /s^2.
}
\end{ex}

% ===== Câu 87 =====
% ID: AIQ_L10C3_FULL_0411
% Mức: TH
\begin{ex}
% ID: AIQ_L10C3_FULL_0411
% Mức: TH
Vật khối lượng $3\,\text{kg}$ chịu hợp lực $9\,\text{N}$. Tính gia tốc (kết quả theo đơn vị m/s^2).
\choice
{4}
{\True 3}
{2}
{6}
\loigiai{
Áp dụng hệ thức phù hợp của bài học, thay số được $3\,\text{m}$ /s^2.
}
\end{ex}

% ===== Câu 88 =====
% ID: AIQ_L10C3_FULL_0412
% Mức: TH
\begin{ex}
% ID: AIQ_L10C3_FULL_0412
% Mức: TH
Vật khối lượng $4\,\text{kg}$ chịu hợp lực $16\,\text{N}$. Tính gia tốc (kết quả theo đơn vị m/s^2).
\choice
{3}
{5}
{\True 4}
{8}
\loigiai{
Áp dụng hệ thức phù hợp của bài học, thay số được $4\,\text{m}$ /s^2.
}
\end{ex}

% ===== Câu 89 =====
% ID: AIQ_L10C3_FULL_0413
% Mức: TH
\begin{ex}
% ID: AIQ_L10C3_FULL_0413
% Mức: TH
Vật khối lượng $5\,\text{kg}$ chịu hợp lực $5\,\text{N}$. Tính gia tốc (kết quả theo đơn vị m/s^2).
\choice
{2}
{2}
{0}
{\True 1}
\loigiai{
Áp dụng hệ thức phù hợp của bài học, thay số được $1\,\text{m}$ /s^2.
}
\end{ex}

% ===== Câu 90 =====
% ID: AIQ_L10C3_FULL_0414
% Mức: VD
\begin{ex}
% ID: AIQ_L10C3_FULL_0414
% Mức: VD
Vật khối lượng $6\,\text{kg}$ chịu hợp lực $12\,\text{N}$. Tính gia tốc (kết quả theo đơn vị m/s^2).
\choice
{\True 2}
{3}
{1}
{4}
\loigiai{
Áp dụng hệ thức phù hợp của bài học, thay số được $2\,\text{m}$ /s^2.
}
\end{ex}

% ===== Câu 91 =====
% ID: AIQ_L10C3_FULL_0415
% Mức: VD
\begin{ex}
% ID: AIQ_L10C3_FULL_0415
% Mức: VD
Vật khối lượng $2\,\text{kg}$ chịu hợp lực $6\,\text{N}$. Tính gia tốc (kết quả theo đơn vị m/s^2).
\choice
{4}
{\True 3}
{2}
{6}
\loigiai{
Áp dụng hệ thức phù hợp của bài học, thay số được $3\,\text{m}$ /s^2.
}
\end{ex}

% ===== Câu 92 =====
% ID: AIQ_L10C3_FULL_0416
% Mức: TH
\begin{ex}
% ID: AIQ_L10C3_FULL_0416
% Mức: TH
Vật khối lượng $3\,\text{kg}$ chịu hợp lực $12\,\text{N}$. Tính gia tốc (kết quả theo đơn vị m/s^2).
\choiceTF
{\True Giá trị cần tìm là $4\,\text{m}$ /s^2.}
{Giá trị cần tìm là $5\,\text{m}$ /s^2.}
{\True Phải xác định đầy đủ các lực hoặc đại lượng liên quan trước khi tính.}
{Có thể bỏ qua điều kiện áp dụng của công thức.}
\loigiai{
\begin{itemchoice}
\item (Đúng) Giá trị cần tìm là $4\,\text{m}$ /s^2. (Vì): kết quả $4\,\text{m}$ /s^2. b) Sai. c) Đúng. d) Sai vì phải kiểm tra điều kiện áp dụng
\item (Sai) Giá trị cần tìm là $5\,\text{m}$ /s^2.
\item (Đúng) Phải xác định đầy đủ các lực hoặc đại lượng liên quan trước khi tính.
\item (Sai) Có thể bỏ qua điều kiện áp dụng của công thức.
\end{itemchoice}
}
\end{ex}

% ===== Câu 93 =====
% ID: AIQ_L10C3_FULL_0417
% Mức: TH
\begin{ex}
% ID: AIQ_L10C3_FULL_0417
% Mức: TH
Vật khối lượng $4\,\text{kg}$ chịu hợp lực $4\,\text{N}$. Tính gia tốc (kết quả theo đơn vị m/s^2).
\choiceTF
{\True Giá trị cần tìm là $1\,\text{m}$ /s^2.}
{Giá trị cần tìm là $2\,\text{m}$ /s^2.}
{\True Phải xác định đầy đủ các lực hoặc đại lượng liên quan trước khi tính.}
{Có thể bỏ qua điều kiện áp dụng của công thức.}
\loigiai{
\begin{itemchoice}
\item (Đúng) Giá trị cần tìm là $1\,\text{m}$ /s^2. (Vì): kết quả $1\,\text{m}$ /s^2. b) Sai. c) Đúng. d) Sai vì phải kiểm tra điều kiện áp dụng
\item (Sai) Giá trị cần tìm là $2\,\text{m}$ /s^2.
\item (Đúng) Phải xác định đầy đủ các lực hoặc đại lượng liên quan trước khi tính.
\item (Sai) Có thể bỏ qua điều kiện áp dụng của công thức.
\end{itemchoice}
}
\end{ex}

% ===== Câu 94 =====
% ID: AIQ_L10C3_FULL_0418
% Mức: VD
\begin{ex}
% ID: AIQ_L10C3_FULL_0418
% Mức: VD
Vật khối lượng $5\,\text{kg}$ chịu hợp lực $10\,\text{N}$. Tính gia tốc (kết quả theo đơn vị m/s^2).
\choiceTF
{\True Giá trị cần tìm là $2\,\text{m}$ /s^2.}
{Giá trị cần tìm là $3\,\text{m}$ /s^2.}
{\True Phải xác định đầy đủ các lực hoặc đại lượng liên quan trước khi tính.}
{Có thể bỏ qua điều kiện áp dụng của công thức.}
\loigiai{
\begin{itemchoice}
\item (Đúng) Giá trị cần tìm là $2\,\text{m}$ /s^2. (Vì): kết quả $2\,\text{m}$ /s^2. b) Sai. c) Đúng. d) Sai vì phải kiểm tra điều kiện áp dụng
\item (Sai) Giá trị cần tìm là $3\,\text{m}$ /s^2.
\item (Đúng) Phải xác định đầy đủ các lực hoặc đại lượng liên quan trước khi tính.
\item (Sai) Có thể bỏ qua điều kiện áp dụng của công thức.
\end{itemchoice}
}
\end{ex}

% ===== Câu 95 =====
% ID: AIQ_L10C3_FULL_0419
% Mức: VD
\begin{ex}
% ID: AIQ_L10C3_FULL_0419
% Mức: VD
Vật khối lượng $6\,\text{kg}$ chịu hợp lực $18\,\text{N}$. Tính gia tốc (kết quả theo đơn vị m/s^2).
\choiceTF
{\True Giá trị cần tìm là $3\,\text{m}$ /s^2.}
{Giá trị cần tìm là $4\,\text{m}$ /s^2.}
{\True Phải xác định đầy đủ các lực hoặc đại lượng liên quan trước khi tính.}
{Có thể bỏ qua điều kiện áp dụng của công thức.}
\loigiai{
\begin{itemchoice}
\item (Đúng) Giá trị cần tìm là $3\,\text{m}$ /s^2. (Vì): kết quả $3\,\text{m}$ /s^2. b) Sai. c) Đúng. d) Sai vì phải kiểm tra điều kiện áp dụng
\item (Sai) Giá trị cần tìm là $4\,\text{m}$ /s^2.
\item (Đúng) Phải xác định đầy đủ các lực hoặc đại lượng liên quan trước khi tính.
\item (Sai) Có thể bỏ qua điều kiện áp dụng của công thức.
\end{itemchoice}
}
\end{ex}

% ===== Câu 96 =====
% ID: AIQ_L10C3_FULL_0420
% Mức: VD
\begin{ex}
% ID: AIQ_L10C3_FULL_0420
% Mức: VD
Vật khối lượng $2\,\text{kg}$ chịu hợp lực $8\,\text{N}$. Tính gia tốc (kết quả theo đơn vị m/s^2).
\choiceTF
{\True Giá trị cần tìm là $4\,\text{m}$ /s^2.}
{Giá trị cần tìm là $5\,\text{m}$ /s^2.}
{\True Phải xác định đầy đủ các lực hoặc đại lượng liên quan trước khi tính.}
{Có thể bỏ qua điều kiện áp dụng của công thức.}
\loigiai{
\begin{itemchoice}
\item (Đúng) Giá trị cần tìm là $4\,\text{m}$ /s^2. (Vì): kết quả $4\,\text{m}$ /s^2. b) Sai. c) Đúng. d) Sai vì phải kiểm tra điều kiện áp dụng
\item (Sai) Giá trị cần tìm là $5\,\text{m}$ /s^2.
\item (Đúng) Phải xác định đầy đủ các lực hoặc đại lượng liên quan trước khi tính.
\item (Sai) Có thể bỏ qua điều kiện áp dụng của công thức.
\end{itemchoice}
}
\end{ex}

% ===== Câu 97 =====
% ID: AIQ_L10C3_FULL_0421
% Mức: TH
\begin{ex}
% ID: AIQ_L10C3_FULL_0421
% Mức: TH
Vật khối lượng $3\,\text{kg}$ chịu hợp lực $3\,\text{N}$. Tính gia tốc (kết quả theo đơn vị m/s^2).
\shortans{1}
\loigiai{
Thay số vào hệ thức của bài học thu được $1\,\text{m}$ /s^2.
}
\end{ex}

% ===== Câu 98 =====
% ID: AIQ_L10C3_FULL_0422
% Mức: TH
\begin{ex}
% ID: AIQ_L10C3_FULL_0422
% Mức: TH
Vật khối lượng $4\,\text{kg}$ chịu hợp lực $8\,\text{N}$. Tính gia tốc (kết quả theo đơn vị m/s^2).
\shortans{2}
\loigiai{
Thay số vào hệ thức của bài học thu được $2\,\text{m}$ /s^2.
}
\end{ex}

% ===== Câu 99 =====
% ID: AIQ_L10C3_FULL_0423
% Mức: TH
\begin{ex}
% ID: AIQ_L10C3_FULL_0423
% Mức: TH
Vật khối lượng $5\,\text{kg}$ chịu hợp lực $15\,\text{N}$. Tính gia tốc (kết quả theo đơn vị m/s^2).
\shortans{3}
\loigiai{
Thay số vào hệ thức của bài học thu được $3\,\text{m}$ /s^2.
}
\end{ex}

% ===== Câu 100 =====
% ID: AIQ_L10C3_FULL_0424
% Mức: VD
\begin{ex}
% ID: AIQ_L10C3_FULL_0424
% Mức: VD
Vật khối lượng $6\,\text{kg}$ chịu hợp lực $24\,\text{N}$. Tính gia tốc (kết quả theo đơn vị m/s^2).
\shortans{4}
\loigiai{
Thay số vào hệ thức của bài học thu được $4\,\text{m}$ /s^2.
}
\end{ex}

% ===== Câu 101 =====
% ID: AIQ_L10C3_FULL_0425
% Mức: VD
\begin{ex}
% ID: AIQ_L10C3_FULL_0425
% Mức: VD
Vật khối lượng $2\,\text{kg}$ chịu hợp lực $2\,\text{N}$. Tính gia tốc (kết quả theo đơn vị m/s^2).
\shortans{1}
\loigiai{
Thay số vào hệ thức của bài học thu được $1\,\text{m}$ /s^2.
}
\end{ex}

% ===== Câu 102 =====
% ID: AIQ_L10C3_FULL_0426
% Mức: VD
\begin{ex}
% ID: AIQ_L10C3_FULL_0426
% Mức: VD
Vật khối lượng $3\,\text{kg}$ chịu hợp lực $6\,\text{N}$. Tính gia tốc (kết quả theo đơn vị m/s^2).
\shortans{2}
\loigiai{
Thay số vào hệ thức của bài học thu được $2\,\text{m}$ /s^2.
}
\end{ex}

% ===== Câu 103 =====
% ID: AIQ_L10C3_FULL_0427
% Mức: VD
\begin{bt}
% ID: AIQ_L10C3_FULL_0427
% Mức: VD
Vật khối lượng $4\,\text{kg}$ chịu hợp lực $12\,\text{N}$. Tính gia tốc (kết quả theo đơn vị m/s^2).
\loigiai{
Xác định các đại lượng đã cho, chọn hệ thức phù hợp và thay số. Kết quả: $3\,\text{m}$ /s^2.
}
\end{bt}

% ===== Câu 104 =====
% ID: AIQ_L10C3_FULL_0428
% Mức: VD
\begin{bt}
% ID: AIQ_L10C3_FULL_0428
% Mức: VD
Vật khối lượng $5\,\text{kg}$ chịu hợp lực $20\,\text{N}$. Tính gia tốc (kết quả theo đơn vị m/s^2).
\loigiai{
Xác định các đại lượng đã cho, chọn hệ thức phù hợp và thay số. Kết quả: $4\,\text{m}$ /s^2.
}
\end{bt}

% ===== Câu 105 =====
% ID: AIQ_L10C3_FULL_0429
% Mức: VD
\begin{bt}
% ID: AIQ_L10C3_FULL_0429
% Mức: VD
Vật khối lượng $6\,\text{kg}$ chịu hợp lực $6\,\text{N}$. Tính gia tốc (kết quả theo đơn vị m/s^2).
\loigiai{
Xác định các đại lượng đã cho, chọn hệ thức phù hợp và thay số. Kết quả: $1\,\text{m}$ /s^2.
}
\end{bt}

% ===== Câu 106 =====
% ID: AIQ_L10C3_FULL_0430
% Mức: VDC
\begin{bt}
% ID: AIQ_L10C3_FULL_0430
% Mức: VDC
Vật khối lượng $2\,\text{kg}$ chịu hợp lực $4\,\text{N}$. Tính gia tốc (kết quả theo đơn vị m/s^2).
\loigiai{
Xác định các đại lượng đã cho, chọn hệ thức phù hợp và thay số. Kết quả: $2\,\text{m}$ /s^2.
}
\end{bt}

% ===== Câu 107 =====
% ID: AIQ_L10C3_FULL_0431
% Mức: VDC
\begin{bt}
% ID: AIQ_L10C3_FULL_0431
% Mức: VDC
Vật khối lượng $3\,\text{kg}$ chịu hợp lực $9\,\text{N}$. Tính gia tốc (kết quả theo đơn vị m/s^2).
\loigiai{
Xác định các đại lượng đã cho, chọn hệ thức phù hợp và thay số. Kết quả: $3\,\text{m}$ /s^2.
}
\end{bt}

% ===== Câu 108 =====
% ID: AIQ_L10C3_FULL_0432
% Mức: VDC
\begin{bt}
% ID: AIQ_L10C3_FULL_0432
% Mức: VDC
Vật khối lượng $4\,\text{kg}$ chịu hợp lực $16\,\text{N}$. Tính gia tốc (kết quả theo đơn vị m/s^2).
\loigiai{
Xác định các đại lượng đã cho, chọn hệ thức phù hợp và thay số. Kết quả: $4\,\text{m}$ /s^2.
}
\end{bt}

% ===== Câu 109 =====
% ID: AIQ_L10C3_FULL_0433
% Mức: NB
\dangbt{Bài toán hệ vật và lực liên kết}
\begin{ex}
% ID: AIQ_L10C3_FULL_0433
% Mức: NB
Vật khối lượng $5\,\text{kg}$ chịu hợp lực $5\,\text{N}$. Tính gia tốc (kết quả theo đơn vị m/s^2).
\choice
{\True 1}
{2}
{0}
{2}
\loigiai{
Áp dụng hệ thức phù hợp của bài học, thay số được $1\,\text{m}$ /s^2.
}
\end{ex}

% ===== Câu 110 =====
% ID: AIQ_L10C3_FULL_0434
% Mức: NB
\begin{ex}
% ID: AIQ_L10C3_FULL_0434
% Mức: NB
Vật khối lượng $6\,\text{kg}$ chịu hợp lực $12\,\text{N}$. Tính gia tốc (kết quả theo đơn vị m/s^2).
\choice
{3}
{\True 2}
{1}
{4}
\loigiai{
Áp dụng hệ thức phù hợp của bài học, thay số được $2\,\text{m}$ /s^2.
}
\end{ex}

% ===== Câu 111 =====
% ID: AIQ_L10C3_FULL_0435
% Mức: NB
\begin{ex}
% ID: AIQ_L10C3_FULL_0435
% Mức: NB
Vật khối lượng $2\,\text{kg}$ chịu hợp lực $6\,\text{N}$. Tính gia tốc (kết quả theo đơn vị m/s^2).
\choice
{2}
{4}
{\True 3}
{6}
\loigiai{
Áp dụng hệ thức phù hợp của bài học, thay số được $3\,\text{m}$ /s^2.
}
\end{ex}

% ===== Câu 112 =====
% ID: AIQ_L10C3_FULL_0436
% Mức: NB
\begin{ex}
% ID: AIQ_L10C3_FULL_0436
% Mức: NB
Vật khối lượng $3\,\text{kg}$ chịu hợp lực $12\,\text{N}$. Tính gia tốc (kết quả theo đơn vị m/s^2).
\choice
{8}
{5}
{3}
{\True 4}
\loigiai{
Áp dụng hệ thức phù hợp của bài học, thay số được $4\,\text{m}$ /s^2.
}
\end{ex}

% ===== Câu 113 =====
% ID: AIQ_L10C3_FULL_0437
% Mức: TH
\begin{ex}
% ID: AIQ_L10C3_FULL_0437
% Mức: TH
Vật khối lượng $4\,\text{kg}$ chịu hợp lực $4\,\text{N}$. Tính gia tốc (kết quả theo đơn vị m/s^2).
\choice
{\True 1}
{2}
{0}
{2}
\loigiai{
Áp dụng hệ thức phù hợp của bài học, thay số được $1\,\text{m}$ /s^2.
}
\end{ex}

% ===== Câu 114 =====
% ID: AIQ_L10C3_FULL_0438
% Mức: TH
\begin{ex}
% ID: AIQ_L10C3_FULL_0438
% Mức: TH
Vật khối lượng $5\,\text{kg}$ chịu hợp lực $10\,\text{N}$. Tính gia tốc (kết quả theo đơn vị m/s^2).
\choice
{3}
{\True 2}
{1}
{4}
\loigiai{
Áp dụng hệ thức phù hợp của bài học, thay số được $2\,\text{m}$ /s^2.
}
\end{ex}

% ===== Câu 115 =====
% ID: AIQ_L10C3_FULL_0439
% Mức: TH
\begin{ex}
% ID: AIQ_L10C3_FULL_0439
% Mức: TH
Vật khối lượng $6\,\text{kg}$ chịu hợp lực $18\,\text{N}$. Tính gia tốc (kết quả theo đơn vị m/s^2).
\choice
{2}
{4}
{\True 3}
{6}
\loigiai{
Áp dụng hệ thức phù hợp của bài học, thay số được $3\,\text{m}$ /s^2.
}
\end{ex}

% ===== Câu 116 =====
% ID: AIQ_L10C3_FULL_0440
% Mức: TH
\begin{ex}
% ID: AIQ_L10C3_FULL_0440
% Mức: TH
Vật khối lượng $2\,\text{kg}$ chịu hợp lực $8\,\text{N}$. Tính gia tốc (kết quả theo đơn vị m/s^2).
\choice
{8}
{5}
{3}
{\True 4}
\loigiai{
Áp dụng hệ thức phù hợp của bài học, thay số được $4\,\text{m}$ /s^2.
}
\end{ex}

% ===== Câu 117 =====
% ID: AIQ_L10C3_FULL_0441
% Mức: VD
\begin{ex}
% ID: AIQ_L10C3_FULL_0441
% Mức: VD
Vật khối lượng $3\,\text{kg}$ chịu hợp lực $3\,\text{N}$. Tính gia tốc (kết quả theo đơn vị m/s^2).
\choice
{\True 1}
{2}
{0}
{2}
\loigiai{
Áp dụng hệ thức phù hợp của bài học, thay số được $1\,\text{m}$ /s^2.
}
\end{ex}

% ===== Câu 118 =====
% ID: AIQ_L10C3_FULL_0442
% Mức: VD
\begin{ex}
% ID: AIQ_L10C3_FULL_0442
% Mức: VD
Vật khối lượng $4\,\text{kg}$ chịu hợp lực $8\,\text{N}$. Tính gia tốc (kết quả theo đơn vị m/s^2).
\choice
{3}
{\True 2}
{1}
{4}
\loigiai{
Áp dụng hệ thức phù hợp của bài học, thay số được $2\,\text{m}$ /s^2.
}
\end{ex}

% ===== Câu 119 =====
% ID: AIQ_L10C3_FULL_0443
% Mức: TH
\begin{ex}
% ID: AIQ_L10C3_FULL_0443
% Mức: TH
Vật khối lượng $5\,\text{kg}$ chịu hợp lực $15\,\text{N}$. Tính gia tốc (kết quả theo đơn vị m/s^2).
\choiceTF
{\True Giá trị cần tìm là $3\,\text{m}$ /s^2.}
{Giá trị cần tìm là $4\,\text{m}$ /s^2.}
{\True Phải xác định đầy đủ các lực hoặc đại lượng liên quan trước khi tính.}
{Có thể bỏ qua điều kiện áp dụng của công thức.}
\loigiai{
\begin{itemchoice}
\item (Đúng) Giá trị cần tìm là $3\,\text{m}$ /s^2. (Vì): kết quả $3\,\text{m}$ /s^2. b) Sai. c) Đúng. d) Sai vì phải kiểm tra điều kiện áp dụng
\item (Sai) Giá trị cần tìm là $4\,\text{m}$ /s^2.
\item (Đúng) Phải xác định đầy đủ các lực hoặc đại lượng liên quan trước khi tính.
\item (Sai) Có thể bỏ qua điều kiện áp dụng của công thức.
\end{itemchoice}
}
\end{ex}

% ===== Câu 120 =====
% ID: AIQ_L10C3_FULL_0444
% Mức: TH
\begin{ex}
% ID: AIQ_L10C3_FULL_0444
% Mức: TH
Vật khối lượng $6\,\text{kg}$ chịu hợp lực $24\,\text{N}$. Tính gia tốc (kết quả theo đơn vị m/s^2).
\choiceTF
{\True Giá trị cần tìm là $4\,\text{m}$ /s^2.}
{Giá trị cần tìm là $5\,\text{m}$ /s^2.}
{\True Phải xác định đầy đủ các lực hoặc đại lượng liên quan trước khi tính.}
{Có thể bỏ qua điều kiện áp dụng của công thức.}
\loigiai{
\begin{itemchoice}
\item (Đúng) Giá trị cần tìm là $4\,\text{m}$ /s^2. (Vì): kết quả $4\,\text{m}$ /s^2. b) Sai. c) Đúng. d) Sai vì phải kiểm tra điều kiện áp dụng
\item (Sai) Giá trị cần tìm là $5\,\text{m}$ /s^2.
\item (Đúng) Phải xác định đầy đủ các lực hoặc đại lượng liên quan trước khi tính.
\item (Sai) Có thể bỏ qua điều kiện áp dụng của công thức.
\end{itemchoice}
}
\end{ex}

% ===== Câu 121 =====
% ID: AIQ_L10C3_FULL_0445
% Mức: VD
\begin{ex}
% ID: AIQ_L10C3_FULL_0445
% Mức: VD
Vật khối lượng $2\,\text{kg}$ chịu hợp lực $2\,\text{N}$. Tính gia tốc (kết quả theo đơn vị m/s^2).
\choiceTF
{\True Giá trị cần tìm là $1\,\text{m}$ /s^2.}
{Giá trị cần tìm là $2\,\text{m}$ /s^2.}
{\True Phải xác định đầy đủ các lực hoặc đại lượng liên quan trước khi tính.}
{Có thể bỏ qua điều kiện áp dụng của công thức.}
\loigiai{
\begin{itemchoice}
\item (Đúng) Giá trị cần tìm là $1\,\text{m}$ /s^2. (Vì): kết quả $1\,\text{m}$ /s^2. b) Sai. c) Đúng. d) Sai vì phải kiểm tra điều kiện áp dụng
\item (Sai) Giá trị cần tìm là $2\,\text{m}$ /s^2.
\item (Đúng) Phải xác định đầy đủ các lực hoặc đại lượng liên quan trước khi tính.
\item (Sai) Có thể bỏ qua điều kiện áp dụng của công thức.
\end{itemchoice}
}
\end{ex}

% ===== Câu 122 =====
% ID: AIQ_L10C3_FULL_0446
% Mức: VD
\begin{ex}
% ID: AIQ_L10C3_FULL_0446
% Mức: VD
Vật khối lượng $3\,\text{kg}$ chịu hợp lực $6\,\text{N}$. Tính gia tốc (kết quả theo đơn vị m/s^2).
\choiceTF
{\True Giá trị cần tìm là $2\,\text{m}$ /s^2.}
{Giá trị cần tìm là $3\,\text{m}$ /s^2.}
{\True Phải xác định đầy đủ các lực hoặc đại lượng liên quan trước khi tính.}
{Có thể bỏ qua điều kiện áp dụng của công thức.}
\loigiai{
\begin{itemchoice}
\item (Đúng) Giá trị cần tìm là $2\,\text{m}$ /s^2. (Vì): kết quả $2\,\text{m}$ /s^2. b) Sai. c) Đúng. d) Sai vì phải kiểm tra điều kiện áp dụng
\item (Sai) Giá trị cần tìm là $3\,\text{m}$ /s^2.
\item (Đúng) Phải xác định đầy đủ các lực hoặc đại lượng liên quan trước khi tính.
\item (Sai) Có thể bỏ qua điều kiện áp dụng của công thức.
\end{itemchoice}
}
\end{ex}

% ===== Câu 123 =====
% ID: AIQ_L10C3_FULL_0447
% Mức: VD
\begin{ex}
% ID: AIQ_L10C3_FULL_0447
% Mức: VD
Vật khối lượng $4\,\text{kg}$ chịu hợp lực $12\,\text{N}$. Tính gia tốc (kết quả theo đơn vị m/s^2).
\choiceTF
{\True Giá trị cần tìm là $3\,\text{m}$ /s^2.}
{Giá trị cần tìm là $4\,\text{m}$ /s^2.}
{\True Phải xác định đầy đủ các lực hoặc đại lượng liên quan trước khi tính.}
{Có thể bỏ qua điều kiện áp dụng của công thức.}
\loigiai{
\begin{itemchoice}
\item (Đúng) Giá trị cần tìm là $3\,\text{m}$ /s^2. (Vì): kết quả $3\,\text{m}$ /s^2. b) Sai. c) Đúng. d) Sai vì phải kiểm tra điều kiện áp dụng
\item (Sai) Giá trị cần tìm là $4\,\text{m}$ /s^2.
\item (Đúng) Phải xác định đầy đủ các lực hoặc đại lượng liên quan trước khi tính.
\item (Sai) Có thể bỏ qua điều kiện áp dụng của công thức.
\end{itemchoice}
}
\end{ex}

% ===== Câu 124 =====
% ID: AIQ_L10C3_FULL_0448
% Mức: TH
\begin{ex}
% ID: AIQ_L10C3_FULL_0448
% Mức: TH
Vật khối lượng $5\,\text{kg}$ chịu hợp lực $20\,\text{N}$. Tính gia tốc (kết quả theo đơn vị m/s^2).
\shortans{4}
\loigiai{
Thay số vào hệ thức của bài học thu được $4\,\text{m}$ /s^2.
}
\end{ex}

% ===== Câu 125 =====
% ID: AIQ_L10C3_FULL_0449
% Mức: TH
\begin{ex}
% ID: AIQ_L10C3_FULL_0449
% Mức: TH
Vật khối lượng $6\,\text{kg}$ chịu hợp lực $6\,\text{N}$. Tính gia tốc (kết quả theo đơn vị m/s^2).
\shortans{1}
\loigiai{
Thay số vào hệ thức của bài học thu được $1\,\text{m}$ /s^2.
}
\end{ex}

% ===== Câu 126 =====
% ID: AIQ_L10C3_FULL_0450
% Mức: TH
\begin{ex}
% ID: AIQ_L10C3_FULL_0450
% Mức: TH
Vật khối lượng $2\,\text{kg}$ chịu hợp lực $4\,\text{N}$. Tính gia tốc (kết quả theo đơn vị m/s^2).
\shortans{2}
\loigiai{
Thay số vào hệ thức của bài học thu được $2\,\text{m}$ /s^2.
}
\end{ex}

% ===== Câu 127 =====
% ID: AIQ_L10C3_FULL_0451
% Mức: VD
\begin{ex}
% ID: AIQ_L10C3_FULL_0451
% Mức: VD
Vật khối lượng $3\,\text{kg}$ chịu hợp lực $9\,\text{N}$. Tính gia tốc (kết quả theo đơn vị m/s^2).
\shortans{3}
\loigiai{
Thay số vào hệ thức của bài học thu được $3\,\text{m}$ /s^2.
}
\end{ex}

% ===== Câu 128 =====
% ID: AIQ_L10C3_FULL_0452
% Mức: VD
\begin{ex}
% ID: AIQ_L10C3_FULL_0452
% Mức: VD
Vật khối lượng $4\,\text{kg}$ chịu hợp lực $16\,\text{N}$. Tính gia tốc (kết quả theo đơn vị m/s^2).
\shortans{4}
\loigiai{
Thay số vào hệ thức của bài học thu được $4\,\text{m}$ /s^2.
}
\end{ex}

% ===== Câu 129 =====
% ID: AIQ_L10C3_FULL_0453
% Mức: VD
\begin{ex}
% ID: AIQ_L10C3_FULL_0453
% Mức: VD
Vật khối lượng $5\,\text{kg}$ chịu hợp lực $5\,\text{N}$. Tính gia tốc (kết quả theo đơn vị m/s^2).
\shortans{1}
\loigiai{
Thay số vào hệ thức của bài học thu được $1\,\text{m}$ /s^2.
}
\end{ex}

% ===== Câu 130 =====
% ID: AIQ_L10C3_FULL_0454
% Mức: VD
\begin{bt}
% ID: AIQ_L10C3_FULL_0454
% Mức: VD
Vật khối lượng $6\,\text{kg}$ chịu hợp lực $12\,\text{N}$. Tính gia tốc (kết quả theo đơn vị m/s^2).
\loigiai{
Xác định các đại lượng đã cho, chọn hệ thức phù hợp và thay số. Kết quả: $2\,\text{m}$ /s^2.
}
\end{bt}

% ===== Câu 131 =====
% ID: AIQ_L10C3_FULL_0455
% Mức: VD
\begin{bt}
% ID: AIQ_L10C3_FULL_0455
% Mức: VD
Vật khối lượng $2\,\text{kg}$ chịu hợp lực $6\,\text{N}$. Tính gia tốc (kết quả theo đơn vị m/s^2).
\loigiai{
Xác định các đại lượng đã cho, chọn hệ thức phù hợp và thay số. Kết quả: $3\,\text{m}$ /s^2.
}
\end{bt}

% ===== Câu 132 =====
% ID: AIQ_L10C3_FULL_0456
% Mức: VD
\begin{bt}
% ID: AIQ_L10C3_FULL_0456
% Mức: VD
Vật khối lượng $3\,\text{kg}$ chịu hợp lực $12\,\text{N}$. Tính gia tốc (kết quả theo đơn vị m/s^2).
\loigiai{
Xác định các đại lượng đã cho, chọn hệ thức phù hợp và thay số. Kết quả: $4\,\text{m}$ /s^2.
}
\end{bt}

% ===== Câu 133 =====
% ID: AIQ_L10C3_FULL_0457
% Mức: VDC
\begin{bt}
% ID: AIQ_L10C3_FULL_0457
% Mức: VDC
Vật khối lượng $4\,\text{kg}$ chịu hợp lực $4\,\text{N}$. Tính gia tốc (kết quả theo đơn vị m/s^2).
\loigiai{
Xác định các đại lượng đã cho, chọn hệ thức phù hợp và thay số. Kết quả: $1\,\text{m}$ /s^2.
}
\end{bt}

% ===== Câu 134 =====
% ID: AIQ_L10C3_FULL_0458
% Mức: VDC
\begin{bt}
% ID: AIQ_L10C3_FULL_0458
% Mức: VDC
Vật khối lượng $5\,\text{kg}$ chịu hợp lực $10\,\text{N}$. Tính gia tốc (kết quả theo đơn vị m/s^2).
\loigiai{
Xác định các đại lượng đã cho, chọn hệ thức phù hợp và thay số. Kết quả: $2\,\text{m}$ /s^2.
}
\end{bt}

% ===== Câu 135 =====
% ID: AIQ_L10C3_FULL_0459
% Mức: VDC
\begin{bt}
% ID: AIQ_L10C3_FULL_0459
% Mức: VDC
Vật khối lượng $6\,\text{kg}$ chịu hợp lực $18\,\text{N}$. Tính gia tốc (kết quả theo đơn vị m/s^2).
\loigiai{
Xác định các đại lượng đã cho, chọn hệ thức phù hợp và thay số. Kết quả: $3\,\text{m}$ /s^2.
}
\end{bt}

% ===== Câu 136 =====
% ID: AIQ_L10C3_FULL_0460
% Mức: NB
\dangbt{Khai thác đồ thị và bài toán thực tiễn về định luật II Newton}
\begin{ex}
% ID: AIQ_L10C3_FULL_0460
% Mức: NB
Vật khối lượng $2\,\text{kg}$ chịu hợp lực $8\,\text{N}$. Tính gia tốc (kết quả theo đơn vị m/s^2).
\choice
{\True 4}
{5}
{3}
{8}
\loigiai{
Áp dụng hệ thức phù hợp của bài học, thay số được $4\,\text{m}$ /s^2.
}
\end{ex}

% ===== Câu 137 =====
% ID: AIQ_L10C3_FULL_0461
% Mức: NB
\begin{ex}
% ID: AIQ_L10C3_FULL_0461
% Mức: NB
Vật khối lượng $3\,\text{kg}$ chịu hợp lực $3\,\text{N}$. Tính gia tốc (kết quả theo đơn vị m/s^2).
\choice
{2}
{\True 1}
{0}
{2}
\loigiai{
Áp dụng hệ thức phù hợp của bài học, thay số được $1\,\text{m}$ /s^2.
}
\end{ex}

% ===== Câu 138 =====
% ID: AIQ_L10C3_FULL_0462
% Mức: NB
\begin{ex}
% ID: AIQ_L10C3_FULL_0462
% Mức: NB
Vật khối lượng $4\,\text{kg}$ chịu hợp lực $8\,\text{N}$. Tính gia tốc (kết quả theo đơn vị m/s^2).
\choice
{1}
{3}
{\True 2}
{4}
\loigiai{
Áp dụng hệ thức phù hợp của bài học, thay số được $2\,\text{m}$ /s^2.
}
\end{ex}

% ===== Câu 139 =====
% ID: AIQ_L10C3_FULL_0463
% Mức: NB
\begin{ex}
% ID: AIQ_L10C3_FULL_0463
% Mức: NB
Vật khối lượng $5\,\text{kg}$ chịu hợp lực $15\,\text{N}$. Tính gia tốc (kết quả theo đơn vị m/s^2).
\choice
{6}
{4}
{2}
{\True 3}
\loigiai{
Áp dụng hệ thức phù hợp của bài học, thay số được $3\,\text{m}$ /s^2.
}
\end{ex}

% ===== Câu 140 =====
% ID: AIQ_L10C3_FULL_0464
% Mức: TH
\begin{ex}
% ID: AIQ_L10C3_FULL_0464
% Mức: TH
Vật khối lượng $6\,\text{kg}$ chịu hợp lực $24\,\text{N}$. Tính gia tốc (kết quả theo đơn vị m/s^2).
\choice
{\True 4}
{5}
{3}
{8}
\loigiai{
Áp dụng hệ thức phù hợp của bài học, thay số được $4\,\text{m}$ /s^2.
}
\end{ex}

% ===== Câu 141 =====
% ID: AIQ_L10C3_FULL_0465
% Mức: TH
\begin{ex}
% ID: AIQ_L10C3_FULL_0465
% Mức: TH
Vật khối lượng $2\,\text{kg}$ chịu hợp lực $2\,\text{N}$. Tính gia tốc (kết quả theo đơn vị m/s^2).
\choice
{2}
{\True 1}
{0}
{2}
\loigiai{
Áp dụng hệ thức phù hợp của bài học, thay số được $1\,\text{m}$ /s^2.
}
\end{ex}

% ===== Câu 142 =====
% ID: AIQ_L10C3_FULL_0466
% Mức: TH
\begin{ex}
% ID: AIQ_L10C3_FULL_0466
% Mức: TH
Vật khối lượng $3\,\text{kg}$ chịu hợp lực $6\,\text{N}$. Tính gia tốc (kết quả theo đơn vị m/s^2).
\choice
{1}
{3}
{\True 2}
{4}
\loigiai{
Áp dụng hệ thức phù hợp của bài học, thay số được $2\,\text{m}$ /s^2.
}
\end{ex}

% ===== Câu 143 =====
% ID: AIQ_L10C3_FULL_0467
% Mức: TH
\begin{ex}
% ID: AIQ_L10C3_FULL_0467
% Mức: TH
Vật khối lượng $4\,\text{kg}$ chịu hợp lực $12\,\text{N}$. Tính gia tốc (kết quả theo đơn vị m/s^2).
\choice
{6}
{4}
{2}
{\True 3}
\loigiai{
Áp dụng hệ thức phù hợp của bài học, thay số được $3\,\text{m}$ /s^2.
}
\end{ex}

% ===== Câu 144 =====
% ID: AIQ_L10C3_FULL_0468
% Mức: VD
\begin{ex}
% ID: AIQ_L10C3_FULL_0468
% Mức: VD
Vật khối lượng $5\,\text{kg}$ chịu hợp lực $20\,\text{N}$. Tính gia tốc (kết quả theo đơn vị m/s^2).
\choice
{\True 4}
{5}
{3}
{8}
\loigiai{
Áp dụng hệ thức phù hợp của bài học, thay số được $4\,\text{m}$ /s^2.
}
\end{ex}

% ===== Câu 145 =====
% ID: AIQ_L10C3_FULL_0469
% Mức: VD
\begin{ex}
% ID: AIQ_L10C3_FULL_0469
% Mức: VD
Vật khối lượng $6\,\text{kg}$ chịu hợp lực $6\,\text{N}$. Tính gia tốc (kết quả theo đơn vị m/s^2).
\choice
{2}
{\True 1}
{0}
{2}
\loigiai{
Áp dụng hệ thức phù hợp của bài học, thay số được $1\,\text{m}$ /s^2.
}
\end{ex}

% ===== Câu 146 =====
% ID: AIQ_L10C3_FULL_0470
% Mức: TH
\begin{ex}
% ID: AIQ_L10C3_FULL_0470
% Mức: TH
Vật khối lượng $2\,\text{kg}$ chịu hợp lực $4\,\text{N}$. Tính gia tốc (kết quả theo đơn vị m/s^2).
\choiceTF
{\True Giá trị cần tìm là $2\,\text{m}$ /s^2.}
{Giá trị cần tìm là $3\,\text{m}$ /s^2.}
{\True Phải xác định đầy đủ các lực hoặc đại lượng liên quan trước khi tính.}
{Có thể bỏ qua điều kiện áp dụng của công thức.}
\loigiai{
\begin{itemchoice}
\item (Đúng) Giá trị cần tìm là $2\,\text{m}$ /s^2. (Vì): kết quả $2\,\text{m}$ /s^2. b) Sai. c) Đúng. d) Sai vì phải kiểm tra điều kiện áp dụng
\item (Sai) Giá trị cần tìm là $3\,\text{m}$ /s^2.
\item (Đúng) Phải xác định đầy đủ các lực hoặc đại lượng liên quan trước khi tính.
\item (Sai) Có thể bỏ qua điều kiện áp dụng của công thức.
\end{itemchoice}
}
\end{ex}

% ===== Câu 147 =====
% ID: AIQ_L10C3_FULL_0471
% Mức: TH
\begin{ex}
% ID: AIQ_L10C3_FULL_0471
% Mức: TH
Vật khối lượng $3\,\text{kg}$ chịu hợp lực $9\,\text{N}$. Tính gia tốc (kết quả theo đơn vị m/s^2).
\choiceTF
{\True Giá trị cần tìm là $3\,\text{m}$ /s^2.}
{Giá trị cần tìm là $4\,\text{m}$ /s^2.}
{\True Phải xác định đầy đủ các lực hoặc đại lượng liên quan trước khi tính.}
{Có thể bỏ qua điều kiện áp dụng của công thức.}
\loigiai{
\begin{itemchoice}
\item (Đúng) Giá trị cần tìm là $3\,\text{m}$ /s^2. (Vì): kết quả $3\,\text{m}$ /s^2. b) Sai. c) Đúng. d) Sai vì phải kiểm tra điều kiện áp dụng
\item (Sai) Giá trị cần tìm là $4\,\text{m}$ /s^2.
\item (Đúng) Phải xác định đầy đủ các lực hoặc đại lượng liên quan trước khi tính.
\item (Sai) Có thể bỏ qua điều kiện áp dụng của công thức.
\end{itemchoice}
}
\end{ex}

% ===== Câu 148 =====
% ID: AIQ_L10C3_FULL_0472
% Mức: VD
\begin{ex}
% ID: AIQ_L10C3_FULL_0472
% Mức: VD
Vật khối lượng $4\,\text{kg}$ chịu hợp lực $16\,\text{N}$. Tính gia tốc (kết quả theo đơn vị m/s^2).
\choiceTF
{\True Giá trị cần tìm là $4\,\text{m}$ /s^2.}
{Giá trị cần tìm là $5\,\text{m}$ /s^2.}
{\True Phải xác định đầy đủ các lực hoặc đại lượng liên quan trước khi tính.}
{Có thể bỏ qua điều kiện áp dụng của công thức.}
\loigiai{
\begin{itemchoice}
\item (Đúng) Giá trị cần tìm là $4\,\text{m}$ /s^2. (Vì): kết quả $4\,\text{m}$ /s^2. b) Sai. c) Đúng. d) Sai vì phải kiểm tra điều kiện áp dụng
\item (Sai) Giá trị cần tìm là $5\,\text{m}$ /s^2.
\item (Đúng) Phải xác định đầy đủ các lực hoặc đại lượng liên quan trước khi tính.
\item (Sai) Có thể bỏ qua điều kiện áp dụng của công thức.
\end{itemchoice}
}
\end{ex}

% ===== Câu 149 =====
% ID: AIQ_L10C3_FULL_0473
% Mức: VD
\begin{ex}
% ID: AIQ_L10C3_FULL_0473
% Mức: VD
Vật khối lượng $5\,\text{kg}$ chịu hợp lực $5\,\text{N}$. Tính gia tốc (kết quả theo đơn vị m/s^2).
\choiceTF
{\True Giá trị cần tìm là $1\,\text{m}$ /s^2.}
{Giá trị cần tìm là $2\,\text{m}$ /s^2.}
{\True Phải xác định đầy đủ các lực hoặc đại lượng liên quan trước khi tính.}
{Có thể bỏ qua điều kiện áp dụng của công thức.}
\loigiai{
\begin{itemchoice}
\item (Đúng) Giá trị cần tìm là $1\,\text{m}$ /s^2. (Vì): kết quả $1\,\text{m}$ /s^2. b) Sai. c) Đúng. d) Sai vì phải kiểm tra điều kiện áp dụng
\item (Sai) Giá trị cần tìm là $2\,\text{m}$ /s^2.
\item (Đúng) Phải xác định đầy đủ các lực hoặc đại lượng liên quan trước khi tính.
\item (Sai) Có thể bỏ qua điều kiện áp dụng của công thức.
\end{itemchoice}
}
\end{ex}

% ===== Câu 150 =====
% ID: AIQ_L10C3_FULL_0474
% Mức: VD
\begin{ex}
% ID: AIQ_L10C3_FULL_0474
% Mức: VD
Vật khối lượng $6\,\text{kg}$ chịu hợp lực $12\,\text{N}$. Tính gia tốc (kết quả theo đơn vị m/s^2).
\choiceTF
{\True Giá trị cần tìm là $2\,\text{m}$ /s^2.}
{Giá trị cần tìm là $3\,\text{m}$ /s^2.}
{\True Phải xác định đầy đủ các lực hoặc đại lượng liên quan trước khi tính.}
{Có thể bỏ qua điều kiện áp dụng của công thức.}
\loigiai{
\begin{itemchoice}
\item (Đúng) Giá trị cần tìm là $2\,\text{m}$ /s^2. (Vì): kết quả $2\,\text{m}$ /s^2. b) Sai. c) Đúng. d) Sai vì phải kiểm tra điều kiện áp dụng
\item (Sai) Giá trị cần tìm là $3\,\text{m}$ /s^2.
\item (Đúng) Phải xác định đầy đủ các lực hoặc đại lượng liên quan trước khi tính.
\item (Sai) Có thể bỏ qua điều kiện áp dụng của công thức.
\end{itemchoice}
}
\end{ex}

% ===== Câu 151 =====
% ID: AIQ_L10C3_FULL_0475
% Mức: TH
\begin{ex}
% ID: AIQ_L10C3_FULL_0475
% Mức: TH
Vật khối lượng $2\,\text{kg}$ chịu hợp lực $6\,\text{N}$. Tính gia tốc (kết quả theo đơn vị m/s^2).
\shortans{3}
\loigiai{
Thay số vào hệ thức của bài học thu được $3\,\text{m}$ /s^2.
}
\end{ex}

% ===== Câu 152 =====
% ID: AIQ_L10C3_FULL_0476
% Mức: TH
\begin{ex}
% ID: AIQ_L10C3_FULL_0476
% Mức: TH
Vật khối lượng $3\,\text{kg}$ chịu hợp lực $12\,\text{N}$. Tính gia tốc (kết quả theo đơn vị m/s^2).
\shortans{4}
\loigiai{
Thay số vào hệ thức của bài học thu được $4\,\text{m}$ /s^2.
}
\end{ex}

% ===== Câu 153 =====
% ID: AIQ_L10C3_FULL_0477
% Mức: TH
\begin{ex}
% ID: AIQ_L10C3_FULL_0477
% Mức: TH
Vật khối lượng $4\,\text{kg}$ chịu hợp lực $4\,\text{N}$. Tính gia tốc (kết quả theo đơn vị m/s^2).
\shortans{1}
\loigiai{
Thay số vào hệ thức của bài học thu được $1\,\text{m}$ /s^2.
}
\end{ex}

% ===== Câu 154 =====
% ID: AIQ_L10C3_FULL_0478
% Mức: VD
\begin{ex}
% ID: AIQ_L10C3_FULL_0478
% Mức: VD
Vật khối lượng $5\,\text{kg}$ chịu hợp lực $10\,\text{N}$. Tính gia tốc (kết quả theo đơn vị m/s^2).
\shortans{2}
\loigiai{
Thay số vào hệ thức của bài học thu được $2\,\text{m}$ /s^2.
}
\end{ex}

% ===== Câu 155 =====
% ID: AIQ_L10C3_FULL_0479
% Mức: VD
\begin{ex}
% ID: AIQ_L10C3_FULL_0479
% Mức: VD
Vật khối lượng $6\,\text{kg}$ chịu hợp lực $18\,\text{N}$. Tính gia tốc (kết quả theo đơn vị m/s^2).
\shortans{3}
\loigiai{
Thay số vào hệ thức của bài học thu được $3\,\text{m}$ /s^2.
}
\end{ex}

% ===== Câu 156 =====
% ID: AIQ_L10C3_FULL_0480
% Mức: VD
\begin{ex}
% ID: AIQ_L10C3_FULL_0480
% Mức: VD
Vật khối lượng $2\,\text{kg}$ chịu hợp lực $8\,\text{N}$. Tính gia tốc (kết quả theo đơn vị m/s^2).
\shortans{4}
\loigiai{
Thay số vào hệ thức của bài học thu được $4\,\text{m}$ /s^2.
}
\end{ex}

% ===== Câu 157 =====
% ID: AIQ_L10C3_FULL_0481
% Mức: VD
\begin{bt}
% ID: AIQ_L10C3_FULL_0481
% Mức: VD
Vật khối lượng $3\,\text{kg}$ chịu hợp lực $3\,\text{N}$. Tính gia tốc (kết quả theo đơn vị m/s^2).
\loigiai{
Xác định các đại lượng đã cho, chọn hệ thức phù hợp và thay số. Kết quả: $1\,\text{m}$ /s^2.
}
\end{bt}

% ===== Câu 158 =====
% ID: AIQ_L10C3_FULL_0482
% Mức: VD
\begin{bt}
% ID: AIQ_L10C3_FULL_0482
% Mức: VD
Vật khối lượng $4\,\text{kg}$ chịu hợp lực $8\,\text{N}$. Tính gia tốc (kết quả theo đơn vị m/s^2).
\loigiai{
Xác định các đại lượng đã cho, chọn hệ thức phù hợp và thay số. Kết quả: $2\,\text{m}$ /s^2.
}
\end{bt}

% ===== Câu 159 =====
% ID: AIQ_L10C3_FULL_0483
% Mức: VD
\begin{bt}
% ID: AIQ_L10C3_FULL_0483
% Mức: VD
Vật khối lượng $5\,\text{kg}$ chịu hợp lực $15\,\text{N}$. Tính gia tốc (kết quả theo đơn vị m/s^2).
\loigiai{
Xác định các đại lượng đã cho, chọn hệ thức phù hợp và thay số. Kết quả: $3\,\text{m}$ /s^2.
}
\end{bt}

% ===== Câu 160 =====
% ID: AIQ_L10C3_FULL_0484
% Mức: VDC
\begin{bt}
% ID: AIQ_L10C3_FULL_0484
% Mức: VDC
Vật khối lượng $6\,\text{kg}$ chịu hợp lực $24\,\text{N}$. Tính gia tốc (kết quả theo đơn vị m/s^2).
\loigiai{
Xác định các đại lượng đã cho, chọn hệ thức phù hợp và thay số. Kết quả: $4\,\text{m}$ /s^2.
}
\end{bt}

% ===== Câu 161 =====
% ID: AIQ_L10C3_FULL_0485
% Mức: VDC
\begin{bt}
% ID: AIQ_L10C3_FULL_0485
% Mức: VDC
Vật khối lượng $2\,\text{kg}$ chịu hợp lực $2\,\text{N}$. Tính gia tốc (kết quả theo đơn vị m/s^2).
\loigiai{
Xác định các đại lượng đã cho, chọn hệ thức phù hợp và thay số. Kết quả: $1\,\text{m}$ /s^2.
}
\end{bt}

% ===== Câu 162 =====
% ID: AIQ_L10C3_FULL_0486
% Mức: VDC
\begin{bt}
% ID: AIQ_L10C3_FULL_0486
% Mức: VDC
Vật khối lượng $3\,\text{kg}$ chịu hợp lực $6\,\text{N}$. Tính gia tốc (kết quả theo đơn vị m/s^2).
\loigiai{
Xác định các đại lượng đã cho, chọn hệ thức phù hợp và thay số. Kết quả: $2\,\text{m}$ /s^2.
}
\end{bt}
